\documentclass[11pt]{article}

\usepackage[margin=1.05in]{geometry}
\usepackage{amsmath,amssymb,amsthm}
\usepackage{mathtools}
\usepackage{lmodern}
\usepackage[T1]{fontenc}
\usepackage{booktabs}
\usepackage{array}
\usepackage{microtype}
\usepackage{xcolor}
\usepackage[colorlinks=true,linkcolor=blue!45!black,citecolor=blue!45!black,urlcolor=blue!45!black]{hyperref}

% ---------------------------------------------------------------------------
% Notation macros
% ---------------------------------------------------------------------------
\newcommand{\seq}[1]{\langle #1\rangle}
\newcommand{\emptyseq}{\seq{\,}}
\newcommand{\len}[1]{\lvert #1\rvert}
\newcommand{\pre}{\sqsubseteq}                 % prefix order on sequences
\newcommand{\cat}{\mathbin{\cdot}}             % sequence concatenation
\newcommand{\Cat}{\bigodot}                    % iterated concatenation
\newcommand{\Blocks}{\mathcal{B}}
\newcommand{\Snap}{\mathbb{S}}
\newcommand{\Wd}{\mathbb{W}}
\newcommand{\Src}{\mathbb{Q}}
\newcommand{\NoFinal}{\bot}
\newcommand{\Undecl}{\bot}
\newcommand{\PhRow}{\boxdot}                   % placeholder row
\newcommand{\BlRow}{\square}                   % blank row
\newcommand{\OvRow}{\boxplus}                  % overflow-summary row
\newcommand{\Bcell}{\beta}                     % blank cell (0,\BlRow)
\newcommand{\Ocell}{\omega_{\OvRow}}
\newcommand{\Mutable}{\mathsf{Mutable}}
\newcommand{\AppendOnly}{\mathsf{AppendOnly}}
\newcommand{\Absent}{\mathsf{Absent}}
\newcommand{\Queued}{\mathsf{Queued}}
\newcommand{\Active}{\mathsf{Active}}
\newcommand{\Finalized}{\mathsf{Finalized}}
\newcommand{\Committed}{\mathsf{Committed}}
\newcommand{\wide}{\mathsf{wide}}
\newcommand{\narrow}{\mathsf{narrow}}
\newcommand{\Preserve}{\mathsf{Preserve}}
\newcommand{\AppendPol}{\mathsf{Append}}
\newcommand{\Rebuild}{\mathsf{Rebuild}}
\newcommand{\None}{\mathsf{None}}
\newcommand{\srcApp}{\mathsf{Append}}
\newcommand{\srcRet}{\mathsf{Retire}}
\newcommand{\srcRep}{\mathsf{Replay}}
\newcommand{\srcRsz}{\mathsf{Resize}}
\newcommand{\srcFail}{\mathsf{FailedWrite}}
\newcommand{\srcExit}{\mathsf{Exit}}
\newcommand{\stRun}{\mathsf{Running}}
\newcommand{\stGrace}{\mathsf{Graceful}}
\newcommand{\stDetach}{\mathsf{Detach}}
\newcommand{\stWF}{\mathsf{WriteFailure}}
\newcommand{\ren}[1]{\rho_{#1}}
\newcommand{\tg}{\mathrm{tag}}
\newcommand{\ntg}{\mathrm{ntag}}
\newcommand{\Un}{U}
\newcommand{\pres}{\mathrm{pres}}
\newcommand{\vis}{\mathrm{vis}}
\newcommand{\ovf}{\mathrm{ovf}}
\newcommand{\newer}{\mathrm{newer}}
\newcommand{\AOK}{\mathcal{A}}
\newcommand{\Canon}{\kappa}
\newcommand{\Press}{\mathrm{press}}
\newcommand{\Req}{\mathrm{req}}
\newcommand{\PH}{\mathrm{ph}}
\newcommand{\Repl}{\mathrm{rep}}
\newcommand{\RRows}{\mathcal{R}}
\newcommand{\Room}{\mathrm{room}}
\newcommand{\Ledger}{L}
\newcommand{\Native}{P}
\newcommand{\Scr}{Q}
\newcommand{\rmode}{m_{\mathrm{r}}}
\newcommand{\rcur}{\ell_{\mathrm{r}}}
\newcommand{\rend}{u_{\mathrm{r}}}
\newcommand{\rpart}{p_{\mathrm{r}}}
\newcommand{\rprep}{\mathit{prep}}
\newcommand{\rcut}{\mathit{cut}}
\newcommand{\flush}{\mathit{flush}}
\newcommand{\shut}{\mathit{shut}}
\newcommand{\run}{\mathit{run}}
\newcommand{\stopr}{\mathit{stop}}
\newcommand{\created}{\gamma}
\newcommand{\MaxLive}{\Lambda}
\newcommand{\MaxRes}{R_{\max}}
\newcommand{\Lmax}{L_{\max}}
\newcommand{\Fmax}{K_{\max}}
\newcommand{\ps}{\mathrm{pv}}
\newcommand{\pc}{\mathrm{cell}}
\DeclareMathOperator{\Res}{Res}
\DeclareMathOperator{\dm}{dm}
\DeclareMathOperator{\last}{last}
\newcommand{\act}[1]{\text{\normalfont\scshape #1}}
\newcommand{\tla}[1]{\texttt{#1}}
\newcommand{\stt}{\Vdash}
\emergencystretch=1.5em

\newtheorem{theorem}{Theorem}[section]
\newtheorem{lemma}[theorem]{Lemma}
\newtheorem{corollary}[theorem]{Corollary}
\newtheorem{proposition}[theorem]{Proposition}
\theoremstyle{definition}
\newtheorem{definition}[theorem]{Definition}
\newtheorem{assumption}[theorem]{Assumption}
\newtheorem{convention}[theorem]{Convention}
\theoremstyle{remark}
\newtheorem{remark}[theorem]{Remark}

\title{\textbf{Elastic Speculative Slots:}\\
A Formally Verified Protocol for In-Order Retirement and\\
Natural Streaming of Concurrent Blocks in a Terminal Viewport}
\author{}
\date{August 24, 2026}

\begin{document}
\maketitle

\begin{abstract}
Interactive terminal applications increasingly multiplex \emph{speculative}
output blocks---streaming model tokens, concurrent build steps, live
spinners---inside a bounded live viewport, while the terminal's own scrollback
demands an immutable, append-only commitment discipline. We present
\emph{Elastic Speculative Slots}, a rendering protocol that reconciles these
constraints through a decoupled three-layer contract: mutable semantic block
state, an exactly-once logical history ledger, and the physical native
scrollback of the host terminal. The protocol admits blocks in order, presents
them in elastically allocated slots under a strict reservation invariant with
graceful overflow summarization and zero-height degradation, streams
append-only blocks line-by-line under row pressure, retires finalized blocks
in order under pressure or flush, and survives width/height resizes via three
policies ($\Preserve$, $\AppendPol$, $\Rebuild$) with a bottom-first,
single-write synchronous replay that is tearing-free by construction. Partial
terminal write failures are fail-stop: the native scrollback receives a prefix
of immutable rows and the commit frontier never advances, precluding
duplicated or reordered retirement. We formalize the protocol as a guarded
transition system, prove a global inductive safety theorem (exact ledger
content, capacity exactness, overflow forcing, final-snapshot immutability,
epoch-scoped native monotonicity, resize logical-neutrality, fail-stop
safety), prove physical realizability on a standard scrolling terminal via a
generic streaming lemma, and establish conditional progress theorems under
weak fairness, including a precise delineation of when pressure-driven
retirement is guaranteed. Every definition and transition corresponds
one-to-one to the mechanically checked TLA\textsuperscript{+} specification
\texttt{ElasticSlots.tla}; TLC exhaustively verified ten invariants and ten
temporal properties over three boundary configurations totaling
$272{,}860$ distinct states with no violations.
\end{abstract}

\tableofcontents

% ===========================================================================
\section{Introduction}\label{sec:intro}
% ===========================================================================

A terminal emulator is a stateful streaming device. Its scrollback is a
line-buffered, append-only log: once a row scrolls off the top of the visible
screen it is frozen in the user's history and cannot be revised. Modern
command-line applications, however, produce output that is intrinsically
\emph{speculative} while it is being produced---a language-model response
being streamed token by token, a package build whose progress bar rewrites
itself, a test runner whose live summary shrinks and grows. Such applications
must therefore maintain a bounded \emph{live viewport} in which content may be
repainted arbitrarily, and a disciplined boundary across which content is
\emph{retired} into scrollback exactly once, in order, and never again
touched.

This paper gives a complete formal treatment of one such discipline, the
\emph{Elastic Speculative Slots} protocol, exactly as specified in the
TLA\textsuperscript{+} module \texttt{ElasticSlots.tla}. The protocol rests on
a decoupled three-layer contract:

\begin{enumerate}
\item \emph{Semantic application state.} A monotonically ordered family of
  blocks $1,\dots,N$, each carrying a speculative snapshot (its \emph{want})
  and, upon finalization, an immutable \emph{final} snapshot.
\item \emph{Logical history ledger.} A width-independent, exactly-once,
  in-order sequence of committed semantic rows, extended only at its end, plus
  the stable emitted prefix of at most one active append-only head block.
\item \emph{Physical native history.} The rows actually scrolled into the
  terminal's scrollback, each tagged with the \emph{proof source} that emitted
  it ($\srcApp$, $\srcRet$, $\srcRep$, $\srcRsz$, $\srcFail$, $\srcExit$) and
  the width at which it was rendered.
\end{enumerate}

The layers are related by invariants, not by identity: the ledger is a pure
function of the commit frontier and the final snapshots; the native history is
an epoch-scoped, append-only physical realization whose rows are always
attributable to committed content, to the stable head prefix, or to explicitly
tagged resize/exit/failure events. Resizes may re-render the entire native
display ($\Rebuild$) without perturbing the ledger by even one row.

\paragraph{Contributions.}
\begin{enumerate}
\item A precise mathematical model of the protocol state, its live-viewport
  geometry (elastic slot allocation, reservation dominance, shrink bridges,
  overflow summarization, zero-height degradation), and its dual history
  ledgers (\S\ref{sec:model}).
\item A guarded-command transition system covering the full lifecycle:
  admission, speculation, allocation, natural append-only streaming,
  finalization, pressure- and flush-driven retirement, three resize policies
  with bottom-first synchronous replay, write-failure fail-stop, and graceful
  and detached shutdown (\S\ref{sec:protocol}), together with a conformance
  table of required and forbidden implementation effects.
\item Full proofs of: exact capacity and overflow forcing
  (Theorem~\ref{thm:capacity}); reservation dominance and bridge termination
  (Lemma~\ref{lem:bridge}); stable-head emission and ledger monotonicity
  (Lemma~\ref{lem:head}); a generic streaming-realization lemma connecting
  abstract native appends to concrete VT scroll mechanics
  (Lemma~\ref{lem:stream}); the bottom-first synchronous replay invariant
  (Lemma~\ref{lem:replay}); a global inductive safety theorem
  (Theorem~\ref{thm:safety}); and conditional progress theorems under weak
  fairness (\S\ref{sec:liveness}), including an explicit counterexample
  delimiting when pure row pressure does \emph{not} force retirement.
\item A one-to-one mapping from the mathematics to the mechanically checked
  actions, invariants, and temporal properties of \texttt{ElasticSlots.tla},
  with exhaustive TLC model-checking results over three boundary
  configurations (\S\ref{sec:tlc}).
\end{enumerate}

% ===========================================================================
\section{System Model and Formal Definitions}\label{sec:model}
% ===========================================================================

\subsection{Preliminaries}

$\mathbb{N}$ denotes the naturals including $0$; for $k\in\mathbb{N}$ we write
$[k]=\{1,\dots,k\}$ (so $[0]=\emptyset$). For a set $X$, $X^{*}$ is the set of
finite sequences over $X$; $\emptyseq$ is the empty sequence; $\len{s}$ is the
length of $s$; $s[j]$ is the $j$-th element ($1$-indexed); and
\[
  s[l..u] \;=\; \seq{s[l],s[l{+}1],\dots,s[u]},
  \qquad\text{with } s[l..u]=\emptyseq \text{ whenever } l>u .
\]
Concatenation is written $s\cat t$, iterated concatenation
$\Cat_{j=l}^{u} s_j = s_l \cat s_{l+1}\cat\cdots\cat s_u$ (empty when $l>u$),
and $x^{k}$ denotes the sequence consisting of $k$ copies of $x$. The
\emph{prefix order} is
\[
  s \pre t \;\iff\; \len{s}\le\len{t} \;\wedge\; \forall j\in[\len{s}]:\,
  s[j]=t[j].
\]
$\last(s)=s[\len{s}]$ for nonempty $s$. All proofs use only these elementary
operations.

\subsection{Parameters and alphabets}

\begin{assumption}[Protocol parameters]\label{ass:params}
The protocol is parameterized by
\begin{itemize}
\item $N\in\mathbb{N}\setminus\{0\}$, the number of block identities,
  $\Blocks=[N]$;
\item $H\in\mathbb{N}\setminus\{0\}$, the maximum transcript (viewport)
  height;
\item $\MaxRes\in\mathbb{N}$, a bound on the number of resize events;
\item $\MaxLive\in\mathbb{N}\setminus\{0\}$, the uncommitted-block pressure
  threshold;
\item a finite nonempty \emph{row alphabet} $\Sigma$ and a finite
  \emph{snapshot universe} $\Snap\subseteq\Sigma^{*}$ with
  $\emptyseq\in\Snap$, containing at least one snapshot of length $1$ and one
  of length $>1$;
\item pairwise-distinct symbols
  $\NoFinal,\PhRow,\BlRow,\OvRow\notin\Sigma$ (respectively: ``no final
  snapshot'', the placeholder row, the blank row, and the overflow-summary
  row).
\end{itemize}
We write $\Lmax=\max\{\len{s}:s\in\Snap\}$ and
$\Fmax = 2N\Lmax$, an upper bound on the physical length of any single write
batch (the factor $2$ arises from narrow-width rendering,
Definition~\ref{def:render}).
\end{assumption}

\begin{definition}[Widths and rendering]\label{def:render}
The width domain is $\Wd=\{\wide,\narrow\}$. The \emph{width rendering
function} $\ren{\omega}:\Sigma^{*}\to\Sigma^{*}$ is
\[
  \ren{\wide}(s)=s,
  \qquad
  \ren{\narrow}(\emptyseq)=\emptyseq,
  \qquad
  \ren{\narrow}(\seq{x}\cat s)=\seq{x,x}\cat\ren{\narrow}(s).
\]
\end{definition}

$\ren{\narrow}$ doubles every semantic row: it is the two-point abstraction of
soft-wrap reflow, in which a semantic row may occupy a different number of
physical terminal rows at different widths. All that the proofs use is that
$\ren{\omega}$ is prefix-monotone and length-homogeneous:

\begin{proposition}\label{prop:render}
For all $s,t\in\Sigma^{*}$ and $\omega\in\Wd$:
(i)~$\len{\ren{\omega}(s)}=\delta_\omega\len{s}$ with $\delta_\wide=1$,
$\delta_\narrow=2$;
(ii)~$s\pre t \Rightarrow \ren{\omega}(s)\pre\ren{\omega}(t)$;
(iii)~$\ren{\omega}(s\cat t)=\ren{\omega}(s)\cat\ren{\omega}(t)$.
\end{proposition}
\begin{proof}
Immediate induction on $\len{s}$ from Definition~\ref{def:render}.
\end{proof}

\begin{definition}[Cells, ledger rows, native rows]\label{def:cells}
Let $\Gamma=\Sigma\cup\{\PhRow,\BlRow,\OvRow\}$ be the \emph{cell alphabet}
and $\Src=\{\srcApp,\srcRet,\srcRep,\srcRsz,\srcFail,\srcExit\}$ the
\emph{source alphabet}. Then:
\begin{itemize}
\item a \emph{cell} is a pair $(o,x)\in(\{0\}\cup\Blocks)\times\Gamma$
  (owner~$0$ marks non-block chrome); the \emph{blank cell} is
  $\Bcell=(0,\BlRow)$ and the \emph{overflow cell} is $\Ocell=(0,\OvRow)$;
\item a \emph{ledger row} is a pair $(i,x)\in\Blocks\times\Sigma$;
\item a \emph{native row} is a quadruple
  $(q,o,x,\omega)\in\Src\times(\{0\}\cup\Blocks)\times\Gamma\times\Wd$.
\end{itemize}
For $i\in\Blocks$ and $s\in\Sigma^{*}$ define the tagging operators
\[
  \tg_i(s) = \seq{(i,s[1]),\dots,(i,s[\len{s}])},
\]
and, with $r=\ren{\omega}(s)$,
\[
  \ntg^{q}_{\omega}(i,s)
    = \seq{(q,i,r[1],\omega),\dots,(q,i,r[\len{r}],\omega)},
\]
and for a cell sequence $u$ the lift
$\ntg^{q}_{\omega}(u)=\seq{(q,o_1,x_1,\omega),\dots}$ where $u[j]=(o_j,x_j)$.
\end{definition}

Note the asymmetry, which is fundamental: \emph{ledger} rows are semantic and
width-independent; \emph{native} rows are physically rendered at a definite
width and carry a provenance source.

\subsection{Protocol state}\label{sec:state}

\begin{definition}[State space]\label{def:state}
A \emph{protocol state} is a tuple with the components of
Table~\ref{tab:state}. We refer to components by their symbols; a primed
symbol denotes the value in the successor state of a transition.
\end{definition}

\begin{table}[htbp]
\centering\footnotesize
\begin{tabular}{@{}lllp{5.9cm}@{}}
\toprule
Symbol & TLA\textsuperscript{+} & Type & Meaning \\
\midrule
$c$ & \tla{c} & $0..N$ & commit frontier: blocks $1..c$ are committed \\
$\varphi$ & \tla{phase} & $\Blocks\to\Phi$ & lifecycle phase;
  $\Phi=\{\Absent,\Queued,\Active,$ $\Finalized,\Committed\}$ \\
$\mu$ & \tla{mode} & $\Blocks\to\mathbb{M}$ &
  presentation contract; $\mathbb{M}=\{\Undecl,\Mutable,$
  $\AppendOnly\}$ \\
$w$ & \tla{want} & $\Blocks\to\Snap$ & speculative snapshot \\
$f$ & \tla{final} & $\Blocks\to\Snap\cup\{\NoFinal\}$ & immutable final
  snapshot \\
$e$ & \tla{emitted} & $\Blocks\to 0..\Lmax$ & stable emitted prefix length
  (append-only head) \\
$\alpha$ & \tla{alloc} & $\Blocks\to 0..H$ & painted slot allocation
  (rows) \\
$\tau$ & \tla{target} & $\Blocks\to 0..H$ & requested slot allocation \\
$\Ledger$ & \tla{history} & $(\Blocks\times\Sigma)^{*}$ & logical history
  ledger \\
$\Native$ & \tla{native} & native rows$^{*}$ & physical native scrollback
  (current epoch) \\
$\omega$ & \tla{width} & $\Wd$ & terminal width class \\
$h$ & \tla{height} & $0..H$ & live transcript height \\
$\nu$ & \tla{resizes} & $0..\MaxRes$ & resize counter \\
$\varepsilon$ & \tla{epoch} & $0..\MaxRes$ & display epoch \\
$\rmode$ & \tla{replayMode} & $\{\None,\AppendPol,\Rebuild\}$ & pending
  replay policy \\
$\rcur,\rend$ & \tla{replayCursor}, \tla{replayEnd} & $0..N{+}1$, $0..N$ &
  replay window $[\rcur..\rend]$ over committed blocks \\
$\rpart$ & \tla{replayPartial} & $0..\Lmax$ & replayed stable-head prefix
  length \\
$\rprep$ & \tla{replayPrepared} & $\mathbb{B}$ & replay frame prepared \\
$\rcut$ & \tla{replayCut} & $0..\Fmax$ & rows that must scroll during replay \\
$\flush$ & \tla{flush} & $\mathbb{B}$ & explicit retirement request \\
$\shut$ & \tla{shutdown} & $\mathbb{B}$ & graceful shutdown initiated \\
$\run$ & \tla{running} & $\mathbb{B}$ & host live \\
$\stopr$ & \tla{stopReason} & $\mathbb{T}$ &
  termination cause; $\mathbb{T}=\{\stRun,\stGrace,$
  $\stDetach,\stWF\}$ \\
\bottomrule
\end{tabular}
\caption{Protocol state components and their counterparts in
\texttt{ElasticSlots.tla}.}
\label{tab:state}
\end{table}

\begin{definition}[Initial state]\label{def:init}
The initial state has $c=0$; $\varphi_i=\Absent$, $\mu_i=\Undecl$,
$w_i=\emptyseq$, $f_i=\NoFinal$, $e_i=\alpha_i=\tau_i=0$ for all $i$;
$\Ledger=\Native=\emptyseq$; $\omega=\wide$; $h=H$; $\nu=\varepsilon=0$;
$\rmode=\None$, $\rcur=\rend=\rpart=\rcut=0$, $\rprep=\mathrm{false}$;
$\flush=\shut=\mathrm{false}$; $\run=\mathrm{true}$; $\stopr=\stRun$.
\end{definition}

\begin{definition}[Derived quantities]\label{def:derived}
In any state:
\begin{align*}
\created &= \bigl|\{\,i\in\Blocks : \varphi_i\ne\Absent\,\}\bigr|
  &&\text{(created count)}\\
\Un_i(s) &= \begin{cases}
  s[\,e_i{+}1\,..\,\len{s}\,] & \text{if } \mu_i=\AppendOnly\\
  s & \text{otherwise}
\end{cases}
  &&\text{(unemitted suffix)}\\
\Repl &\;\iff\; \rmode\ne\None &&\text{(replay pending).}
\end{align*}
\end{definition}

The lifecycle invariant (Theorem~\ref{thm:safety}) guarantees that the created
blocks are exactly $[\created]$ and that they are partitioned as
$1..c$ committed, $c{+}1..\created$ live (queued, active, or finalized),
$\created{+}1..N$ absent. We freely use this shape below.

\subsection{Presentation contracts}

Every created block carries one of two contracts, fixed at creation:

\begin{itemize}
\item $\Mutable$: while $\varphi_i\in\{\Queued,\Active\}$ the snapshot $w_i$
  may be replaced by \emph{any} member of $\Snap$; the block's viewport slot
  is repainted in place. No row of a mutable block may reach scrollback before
  the whole block retires.
\item $\AppendOnly$: the snapshot evolves only by prefix extension
  ($w_i\pre w_i'$). Consequently an ever-growing \emph{stable prefix} exists,
  and while the block is the \emph{head} (the first uncommitted block,
  $i=c{+}1$) its stable rows may scroll naturally, line by line, into
  scrollback ahead of finalization---this is the streaming affordance.
\end{itemize}

\subsection{The two histories}\label{sec:histories}

\begin{definition}[Committed ledger and stable head]\label{def:ledger}
For $k\in 0..N$ define the \emph{committed ledger}
\[
  C(k) \;=\; \Cat_{j=1}^{k}\, \tg_j(f_j).
\]
The \emph{partial-head predicate} is
\[
  \PH \;\iff\;
  c<\created \,\wedge\, \mu_{c+1}=\AppendOnly \,\wedge\,
  \varphi_{c+1}\in\{\Active,\Finalized\} \,\wedge\, e_{c+1}>0,
\]
and the corresponding \emph{active head} prefix is
\[
  A(c) \;=\;
  \begin{cases}
    \tg_{c+1}\bigl(w_{c+1}[1..e_{c+1}]\bigr) & \text{if } \PH\\
    \emptyseq & \text{otherwise.}
  \end{cases}
\]
\end{definition}

The central ledger invariant, proved inductive in
Theorem~\ref{thm:safety}, is the \emph{exactness equation}
\begin{equation}\label{eq:ech}
  \Ledger \;=\; C(c)\cat A(c).
\end{equation}
The logical history is thus never an independent data structure: it is at all
times \emph{exactly} the committed finals in block order followed by the
stable emitted prefix of the single append-only head, with no duplication,
omission, or reordering. The physical history $\Native$, by contrast, is an
epoch-scoped append-only realization: within a display epoch it only grows at
the end; the $\Rebuild$ resize policy resets it to $\emptyseq$ and increments
$\varepsilon$.

\subsection{Live geometry}\label{sec:geometry}

Throughout this subsection fix the ambient tuple
$(\varphi,f,e,\omega,h)$; all geometric operators are functions of it (the
TLA\textsuperscript{+} operators take these as explicit parameters, which is
essential for stating transition guards over \emph{successor} values).

\begin{definition}[Presentation]\label{def:presented}
Block $i$ is \emph{presented}, written $\pres(i)$, iff
\[
  \varphi_i=\Active
  \;\vee\;
  \bigl(\varphi_i=\Finalized \wedge
        \len{\ren{\omega}(\Un_i(f_i))}>0\bigr).
\]
Let $\pi=\bigl|\{i:\pres(i)\}\bigr|$ be the presented count. The viewport
\emph{overflows}, written $\ovf$, iff $\pi>h$; the \emph{summary row count} is
\[
  \sigma \;=\; \begin{cases} 1 & \text{if } h>0 \wedge \ovf\\
                             0 & \text{otherwise.}\end{cases}
\]
With $\newer(i)=\bigl|\{\,j>i : \pres(j)\,\}\bigr|$, block $i$ is
\emph{visible}, written $\vis(i)$, iff
\[
  \pres(i) \;\wedge\;
  \bigl(\ovf \Rightarrow h>0 \wedge \newer(i)<h-1\bigr).
\]
\end{definition}

Under overflow, visibility privileges recency: only the $h-1$ presented blocks
with the largest indices remain visible, and one row is sacrificed to a
summary marker $\OvRow$ standing for everything older.

\begin{definition}[Allocation admissibility]\label{def:aok}
Allocation vectors $\alpha,\tau\in(0..H)^{\Blocks}$ are \emph{admissible},
written $\AOK(\alpha,\tau)$ (with the ambient tuple implicit), iff
\begin{enumerate}
\item for every $i$ with $\vis(i)$: if $\varphi_i=\Active$ then
  $\alpha_i\in 1..H$ and $\tau_i\in 1..H$; otherwise ($\varphi_i=\Finalized$)
  $\alpha_i\in 1..H$ and $\tau_i=\alpha_i$;
\item for every $i$ with $\neg\vis(i)$: $\alpha_i=\tau_i=0$;
\item the \emph{reservation invariant} holds:
\begin{equation}\label{eq:reservation}
  \Res(\alpha,\tau)+\sigma \;\le\; h,
  \qquad\text{where}\quad
  \Res(\alpha,\tau)=\sum_{i=1}^{N}\max(\alpha_i,\tau_i).
\end{equation}
\end{enumerate}
The \emph{canonical allocation} is
$\Canon = \bigl(i\mapsto \text{if } \vis(i) \text{ then } 1
\text{ else } 0\bigr)$.
\end{definition}

The reservation invariant charges every block for the \emph{maximum} of its
painted and requested heights: growth is admitted only if the grown extent is
already reserved, and a shrinking block retains its old reservation until the
shrink is actually painted. This is what makes concurrent elastic animation
safe (Lemma~\ref{lem:bridge}).

\begin{definition}[Row demand and pressure]\label{def:pressure}
The \emph{row demand} of block $i$ and the total demand are
\[
  \dm(i)=\begin{cases}
    \max\bigl(1,\len{\ren{\omega}(\Un_i(w_i))}\bigr)
      & \varphi_i\in\{\Queued,\Active\}\\
    \len{\ren{\omega}(\Un_i(f_i))}
      & \varphi_i=\Finalized\\
    0 & \text{otherwise,}
  \end{cases}
  \qquad
  D=\sum_{i=1}^{N}\dm(i).
\]
Then \emph{row pressure} is $D>h$; \emph{pressure} is
$\Press \iff D>h \vee \created-c\ge\MaxLive$; and \emph{retirement is
requested} iff $\Req \iff \flush\vee\Press$.
\end{definition}

\begin{definition}[Screen layout]\label{def:screen}
The \emph{preview source} and \emph{preview cell} of block $i$ are
\[
  \ps(i)=\begin{cases}\Un_i(w_i)&\varphi_i=\Active\\ \Un_i(f_i)&
  \text{otherwise,}\end{cases}
  \qquad
  \pc(i)=\begin{cases}
    \bigl(i,\;\last(\ren{\omega}(\ps(i)))\bigr)
      & \ren{\omega}(\ps(i))\ne\emptyseq\\
    (i,\PhRow) & \text{otherwise.}
  \end{cases}
\]
With $A_{\Sigma}=\sum_i \alpha_i$, the \emph{screen} is the cell sequence
\begin{equation}\label{eq:screen}
  \Scr \;=\; \Bcell^{\,h-A_{\Sigma}-\sigma}
  \;\cat\; \Ocell^{\,\sigma}
  \;\cat\; \Cat_{i=1}^{N} \pc(i)^{\alpha_i}.
\end{equation}
\end{definition}

The model abstracts a slot's visual content to $\alpha_i$ copies of a
representative cell carrying the owner and the last rendered unemitted row; a
conforming implementation displays the $\alpha_i$-row tail window of
$\ren{\omega}(\ps(i))$ (padded with the placeholder when empty). Blank rows
sit at the top, then the overflow marker if any, then the slots in ascending
block order---the newest block is adjacent to the cursor at the bottom.
Well-definedness of \eqref{eq:screen} (in particular that $\alpha_i>0$ implies
$\ps(i)$ is a genuine snapshot and the exponent of $\Bcell$ is nonnegative) is
part of Theorem~\ref{thm:capacity}.

\subsection{Replay geometry}\label{sec:replaygeo}

\begin{definition}[Replay frame]\label{def:replayrows}
When $\Repl$ holds, the \emph{replay frame} is
\[
  \RRows \;=\;
  \Bigl(\Cat_{j=\rcur}^{\rend} \ntg^{\srcRep}_{\omega}(j,f_j)\Bigr)
  \;\cat\;
  \ntg^{\srcRep}_{\omega}\bigl(\rend{+}1,\; w_{\rend+1}[1..\rpart]\bigr),
\]
(the second factor being empty when $\rpart=0$), and $\RRows=\emptyseq$
otherwise. The \emph{replay room} and \emph{required cut} are
\[
  \Room=\bigl|\{\,j\in[h] : \Scr[j]=\Bcell\,\}\bigr|,
  \qquad
  \rcut^{*}=\max\bigl(0,\;\len{\RRows}-\Room\bigr).
\]
When $\rprep$ holds, the \emph{prepared tail} is
$\RRows[\rcut{+}1\,..\,\len{\RRows}]$.
\end{definition}

The replay frame re-renders, at the \emph{current} width, all committed finals
in the window $[\rcur..\rend]$ (invariantly $\rcur=1$, $\rend\le c$) plus the
already-emitted stable prefix of the head. The cut $\rcut^{*}$ is the exact
number of frame rows that cannot fit in the blank region of the viewport and
must therefore scroll into native history.

% ===========================================================================
\section{The Rendering Protocol as a Transition System}\label{sec:protocol}
% ===========================================================================

\begin{convention}\label{conv:trans}
Each transition below is a guarded command over the state of
Definition~\ref{def:state}. A transition may fire only when its guard holds;
its effect lists the components it changes; \emph{all unmentioned components
are unchanged}. Every guard implicitly conjoins $\run$. Derived operators
($\vis$, $\sigma$, $\Canon$, $\AOK$, $\Scr$, \dots) evaluated in the effect of
a transition are evaluated at the \emph{successor} values of their arguments;
where a guard constrains a successor configuration we make the substitution
explicit. Substitution into a function is written
$\varphi[i\mapsto p]$.
\end{convention}

\subsection{Admission and speculation}

\paragraph{\act{Create}$(m)$, $m\in\{\Mutable,\AppendOnly\}$.}
\emph{Guard:} $\neg\shut \wedge \created<N$.
\emph{Effect:}
\[
  \varphi_{\created+1}':=\Queued,\qquad \mu_{\created+1}':=m.
\]
Blocks are created contiguously in index order; the index order is the
commitment order.

\paragraph{\act{Admit}$(i)$.}
\emph{Guard:} $\neg\shut \wedge \varphi_i=\Queued \wedge
\forall j<i:\varphi_j\ne\Queued$, and with
$\varphi^{\dagger}=\varphi[i\mapsto\Active]$,
$\alpha^{\dagger}=\alpha[i\mapsto 1]$, $\tau^{\dagger}=\tau[i\mapsto 1]$:
\[
  \neg\ovf(\varphi^{\dagger},f,e,\omega,h)
  \;\wedge\;
  \AOK(\alpha^{\dagger},\tau^{\dagger})
  \text{ at } (\varphi^{\dagger},f,e,\omega,h).
\]
\emph{Effect:} $\varphi':=\varphi^{\dagger}$, $\alpha':=\alpha^{\dagger}$,
$\tau':=\tau^{\dagger}$.
Admission is FIFO (no earlier block is still queued) and \emph{denied} rather
than summarized when it would overflow the viewport: overflow may arise only
from finalization of queued blocks or from resize, never from admission.

\paragraph{\act{Update}$(i,s)$, $s\in\Snap$.}
\emph{Guard:} $\neg\shut \wedge \varphi_i\in\{\Queued,\Active\} \wedge
(\mu_i=\Mutable \vee w_i\pre s) \wedge s\ne w_i$.
\emph{Effect:} $w_i':=s$.
This is the only transition by which speculation evolves; the prefix guard is
the entirety of the append-only contract.

\subsection{Elastic allocation}

\paragraph{\act{RequestAllocation}$(\tau^{\dagger})$,
$\tau^{\dagger}\in(0..H)^{\Blocks}$.}
\emph{Guard:} $\neg\shut \wedge \AOK(\alpha,\tau^{\dagger}) \wedge
\tau^{\dagger}\ne\tau$.
\emph{Effect:} $\tau':=\tau^{\dagger}$.
Targets may change only if the joint reservation
\eqref{eq:reservation} of the \emph{current} paint and the \emph{new} request
still fits.

\begin{definition}[Shrink bridge]\label{def:bridge}
The \emph{bridge function} $B:(0..H)^2\to 0..H$ is
\[
  B(a,t)=\begin{cases}
    t & a<t \quad\text{(growth is immediate)}\\
    2 & a>2 \wedge t=1 \quad\text{(deep shrink bridges through 2)}\\
    t & \text{otherwise.}
  \end{cases}
\]
\end{definition}

\paragraph{\act{ApplyAllocation}$(i)$.}
\emph{Guard:} $\neg\shut \wedge \varphi_i=\Active \wedge \alpha_i\ne\tau_i
\wedge \AOK(\alpha[i\mapsto B(\alpha_i,\tau_i)],\tau)$.
\emph{Effect:} $\alpha_i':=B(\alpha_i,\tau_i)$.
A slot collapsing from $k>2$ rows to $1$ paints an intermediate $2$-row frame,
preserving animation stability (the user perceives a contraction, not a
snap); Lemma~\ref{lem:bridge} shows the bridge terminates and never violates
the reservation.

\subsection{Finalization}

\paragraph{\act{FinalizeActive}$(i,s)$, $s\in\Snap$.}
\emph{Guard:} $\neg\shut \wedge \varphi_i=\Active \wedge
(\mu_i=\Mutable \vee w_i\pre s)$.
\emph{Effect:}
\[
  \varphi_i':=\Finalized,\quad w_i':=s,\quad f_i':=s,\quad
  \alpha':=\tau':=\Canon \text{ at } (\varphi',f',e,\omega,h).
\]

\paragraph{\act{FinalizeQueued}$(i,s)$, $s\in\Snap$.}
Identical, with guard $\varphi_i=\Queued$: a queued block may complete without
ever being admitted to a live slot (e.g.\ a task that finished before any
viewport space became available). Finalization freezes the semantic value
($f_i$ is never rewritten; Theorem~\ref{thm:actions}) and collapses \emph{all}
allocations to canonical single-row slots---finalized content no longer
animates.

\subsection{Natural streaming of the append-only head}

\paragraph{\act{AppendStable}.}
\emph{Guard:}
\[
  \neg\shut \wedge \neg\Repl \wedge c<\created \wedge
  \mu_{c+1}=\AppendOnly \wedge \varphi_{c+1}\in\{\Active,\Finalized\}
  \wedge D>h \wedge e_{c+1}<\len{w_{c+1}}.
\]
\emph{Effect:} with $n=e_{c+1}+1$,
\[
  \Ledger' := \Ledger\cat\seq{(c{+}1,\,w_{c+1}[n])},\qquad
  \Native' := \Native\cat\ntg^{\srcApp}_{\omega}
      \bigl(c{+}1,\,w_{c+1}[n..n]\bigr),
\]
\[
  e_{c+1}' := n,\qquad
  \alpha':=\tau':=\Canon \text{ at } (\varphi,f,e',\omega,h).
\]
One stable row of the head scrolls naturally into both histories. The
transition is enabled only under \emph{row} pressure: while the viewport has
room, stable rows stay live (and remain revisable in presentation, though not
in value). Only the head block ever streams: this preserves in-order
retirement.

\paragraph{\act{CompleteAppendOnly}.}
\emph{Guard:} $\neg\Repl \wedge c<\created \wedge \mu_{c+1}=\AppendOnly
\wedge \varphi_{c+1}=\Finalized \wedge e_{c+1}=\len{f_{c+1}}$.
\emph{Effect:}
\[
  c':=c+1,\quad \varphi_{c+1}':=\Committed,\quad e_{c+1}':=0,\quad
  \alpha':=\tau':=\Canon.
\]
When the head has streamed its entire final snapshot, commitment is pure
bookkeeping---no write occurs, because every row is already in both
histories. Note the guard deliberately omits $\neg\shut$: draining the head
must remain possible during graceful shutdown.

\subsection{Retirement}

\paragraph{\act{BeginFlush}.}
\emph{Guard:} $\neg\flush$. \emph{Effect:} $\flush':=\mathrm{true}$. The flag
is never reset.

\begin{definition}[Retirement batch]\label{def:retbatch}
For $b\in(c{+}1)..N$ with $\varphi_j=\Finalized$ for all $j\in c{+}1..b$, the
\emph{logical} and \emph{physical retirement rows} are
\[
  F_{\log}(b) = \tg_{c+1}\bigl(f_{c+1}[\,e_{c+1}{+}1..\len{f_{c+1}}\,]\bigr)
    \cat \Cat_{j=c+2}^{b}\tg_j(f_j),
\]
\[
  F^{q}_{\mathrm{phys}}(b) =
    \ntg^{q}_{\omega}\bigl(c{+}1,\,
      f_{c+1}[\,e_{c+1}{+}1..\len{f_{c+1}}\,]\bigr)
    \cat \Cat_{j=c+2}^{b}\ntg^{q}_{\omega}(j,f_j),
  \qquad q\in\{\srcRet,\srcFail\}.
\]
\end{definition}

The first (head) block contributes only its \emph{unemitted} suffix---rows
already streamed by \act{AppendStable} are never re-emitted.

\paragraph{\act{RetireSuccess}$(b)$.}
\emph{Guard:} $\neg\Repl \wedge b\in(c{+}1)..N \wedge
\bigl(\forall j\in c{+}1..b:\varphi_j=\Finalized\bigr) \wedge \Req$.
\emph{Effect:}
\[
  \Ledger' := \Ledger\cat F_{\log}(b),\qquad
  \Native' := \Native\cat F^{\srcRet}_{\mathrm{phys}}(b),
\]
\[
  c' := b,\qquad
  \varphi_j' := \Committed \text{ and } e_j' := 0 \text{ for } j\le b,\qquad
  \alpha':=\tau':=\Canon .
\]
Retirement advances the frontier across a \emph{contiguous finalized prefix}
and is permitted only under demand ($\Req$): pressure (row demand exceeding
the viewport, or too many uncommitted blocks) or an explicit flush.
Section~\ref{sec:realization} shows the physical effect is realized by a
single streamed write.

\paragraph{\act{RetireFailure}$(b,k)$.}
\emph{Guard:} as \act{RetireSuccess}$(b)$, plus
$k\in 0..\len{F^{\srcFail}_{\mathrm{phys}}(b)}$.
\emph{Effect:}
\[
  \Native' := \Native\cat F^{\srcFail}_{\mathrm{phys}}(b)[1..k],\qquad
  \run':=\mathrm{false},\qquad \stopr':=\stWF .
\]
$\Ledger$, $c$, $\varphi$, $e$ are \emph{unchanged}: a partial terminal write
deposits an arbitrary prefix of immutable finalized/stable rows into native
scrollback, tagged $\srcFail$, and the host halts. Because the frontier did
not advance, no retry can duplicate rows and no later block can overtake the
failed batch---failure is fail-stop, not fail-retry
(Theorem~\ref{thm:actions}).

\subsection{Resize and synchronous bottom-first replay}

\paragraph{\act{Resize}$(\omega^{\dagger},h^{\dagger},p,k)$,
$\omega^{\dagger}\in\Wd$, $h^{\dagger}\in 0..H$,
$p\in\{\Preserve,\AppendPol,\Rebuild\}$, $k\in 0..\len{\Scr}$.}
\emph{Guard:} $\neg\shut \wedge \nu<\MaxRes \wedge
(\omega^{\dagger}\ne\omega \vee h^{\dagger}\ne h)$.
\emph{Effect:} let
\[
  p^{\dagger} = \begin{cases} p & \omega^{\dagger}\ne\omega\\
                              \Preserve & \text{otherwise,}\end{cases}
  \qquad
  \mathit{begin} \iff p^{\dagger}\ne\Preserve \wedge (c>0 \vee \PH),
\]
and set
\begin{align*}
  \omega' &:= \omega^{\dagger}, \quad h' := h^{\dagger}, \quad \nu':=\nu+1,
  \quad \alpha':=\tau':=\Canon \text{ at } (\varphi,f,e,\omega',h'),\\
  \Native' &:= \begin{cases}
      \emptyseq & p^{\dagger}=\Rebuild\\
      \Native\cat\ntg^{\srcRsz}_{\omega}(\Scr[1..k]) & \text{otherwise,}
    \end{cases}
  \qquad
  \varepsilon' := \begin{cases}\varepsilon+1 & p^{\dagger}=\Rebuild\\
                               \varepsilon & \text{otherwise,}\end{cases}\\
  (\rmode',\rcur',\rend',\rpart') &:=
    \begin{cases}
      \bigl(p^{\dagger},\,1,\,c,\;
        \text{if } \PH \text{ then } e_{c+1} \text{ else } 0\bigr)
        & \mathit{begin}\\
      (\rmode,\rcur,\rend,\rpart) & \neg\mathit{begin}\wedge\Repl\\
      (\None,0,0,0) & \text{otherwise,}
    \end{cases}\\
  \rprep' &:= \mathrm{false},\qquad \rcut':=0 .
\end{align*}
Height-only resizes are always $\Preserve$. The parameter $k$ captures
\emph{emulator-induced push}: when the host terminal itself scrolls up to $k$
top viewport rows into its scrollback during the resize (e.g.\ on height
shrink), those rows are recorded in $\Native$ tagged $\srcRsz$---the logical
ledger is untouched (Theorem~\ref{thm:actions},
\emph{resize logical-neutrality}). Under $\AppendPol$, committed history is
re-rendered at the new width and appended; under $\Rebuild$, the display epoch
is reset and native history is reconstructed from scratch. A resize arriving
\emph{during} an unfinished replay preserves the replay window but invalidates
any prepared frame.

\paragraph{\act{PrepareReplay}.}
\emph{Guard:} $\Repl \wedge \neg\rprep$.
\emph{Effect:} $\rprep':=\mathrm{true}$, $\rcut':=\rcut^{*}$.
Preparation samples the current viewport, computes the replay frame $\RRows$,
and splits it at the required cut: the tail
$\RRows[\rcut^{*}{+}1..]$ will be painted \emph{bottom-first} into the blank
region (no scroll, no history effect), and only the prefix
$\RRows[1..\rcut^{*}]$ will scroll.

\paragraph{\act{ReplaySynchronousSuccess}.}
\emph{Guard:} $\Repl \wedge \rprep$.
\emph{Effect:}
\[
  \Native' := \Native\cat\RRows[1..\rcut],\qquad
  (\rmode',\rcur',\rend',\rpart',\rprep',\rcut') :=
  (\None,0,0,0,\mathrm{false},0).
\]

\paragraph{\act{ReplaySynchronousFailure}$(k)$, $k\in 0..\rcut$.}
\emph{Guard:} $\Repl \wedge \rprep$.
\emph{Effect:} $\Native':=\Native\cat\RRows[1..k]$,
$\run':=\mathrm{false}$, $\stopr':=\stWF$.

\subsection{Shutdown}

\paragraph{\act{BeginGracefulShutdown}.}
\emph{Guard:} $\neg\shut$.
\emph{Effect:} for every $i$ with $\varphi_i\ne\Absent$,
\[
  \varphi_i' := \begin{cases}\Committed & i\le c\\ \Finalized &
  \text{otherwise,}\end{cases}
  \qquad
  f_i' := \begin{cases}f_i & i\le c \vee \varphi_i=\Finalized\\
                       w_i & \text{otherwise,}\end{cases}
\]
\[
  \alpha':=\tau':=\Canon \text{ at } (\varphi',f',e,\omega,h),
  \qquad \flush':=\mathrm{true},\qquad \shut':=\mathrm{true}.
\]
All queued and active work is frozen at its current speculation ($f:=w$) and
the pipeline is put under permanent flush; the remaining trajectory can only
drain (stream the head, commit, replay) and exit.

\paragraph{\act{GracefulExit}$(k)$, $k\in\{0,1\}$.}
\emph{Guard:} $\shut \wedge \neg\Repl \wedge c=\created \wedge
(k=0 \vee h>0)$.
\emph{Effect:} $\run':=\mathrm{false}$, $\stopr':=\stGrace$, and if $k=1$,
$\Native':=\Native\cat\ntg^{\srcExit}_{\omega}(\seq{\Scr[1]})$.
Exit is permitted only after every created block has committed and any replay
has drained; the optional single pushed row (tagged $\srcExit$) models the
final scroll that blanks the viewport and restores the shell prompt.

\paragraph{\act{DetachExit}$(k)$, $k\in\{0,1\}$.}
\emph{Guard:} $\neg\shut \wedge (k=0 \vee h>0)$.
\emph{Effect:} as \act{GracefulExit} with $\stopr':=\stDetach$. Detachment
abandons uncommitted work without emitting it; the exactness equation
\eqref{eq:ech} guarantees the ledger holds precisely the committed content at
the instant of detachment.

\subsection{Scheduler gate and fairness}\label{sec:next}

\begin{definition}[Next-state relation and specification]\label{def:spec}
Let $\mathcal{T}$ be the set of all transitions above except the two replay
writes. The next-state relation is the \emph{gated} disjunction
\[
  \mathrm{Next} \;\equiv\;
  \begin{cases}
    \act{ReplaySynchronousSuccess} \vee
    \exists k:\act{ReplaySynchronousFailure}(k)
      & \rprep\\[2pt]
    \bigvee_{T\in\mathcal{T}} \exists\,\text{params}: T
      & \neg\rprep .
  \end{cases}
\]
The specification is
\[
  \begin{aligned}
  \mathrm{Spec} \;\equiv\;
  \mathrm{Init} \wedge \Box[\mathrm{Next}]_{\mathit{vars}}
  &\wedge \mathrm{WF}(\act{RetireSuccess}^{\exists})
  \wedge \mathrm{WF}(\act{PrepareReplay})\\
  &\wedge \mathrm{WF}(\act{ReplaySynchronousSuccess})
  \wedge \mathrm{WF}(\act{AppendStable})\\
  &\wedge \mathrm{WF}(\act{CompleteAppendOnly}),
  \end{aligned}
\]
where $\act{RetireSuccess}^{\exists}=\exists b:\act{RetireSuccess}(b)$ and
$\mathrm{WF}$ is weak fairness: an action that is eventually forever enabled
is taken infinitely often.
\end{definition}

The gate is the formal content of the word \emph{synchronous}: once a replay
frame is prepared, \emph{no} other transition---no update, no allocation, no
repaint-inducing step---may interleave before the single buffered write
completes or fails. Section~\ref{sec:replayproof} shows this gate is
load-bearing: the prepared-frame invariant $\rcut=\rcut^{*}$ would not be
stable without it.

\subsection{Conformance summary}\label{sec:conformance}

Table~\ref{tab:conformance} summarizes, per transition, the physical effects
a conforming implementation \emph{must} produce and those it \emph{must not}
produce. ``Repaint only'' means: the implementation may rewrite the $h$ live
viewport rows in place but must not cause any row to enter scrollback and
must not touch the ledger.

\begin{table}[htbp]
\centering\footnotesize
\begin{tabular}{@{}>{\raggedright\arraybackslash}p{4.35cm}
                 >{\raggedright\arraybackslash}p{4.95cm}
                 >{\raggedright\arraybackslash}p{4.95cm}@{}}
\toprule
Transition & Required effect & Forbidden effect \\
\midrule
\act{Create}, \act{Update},
\act{RequestAllocation}, \act{BeginFlush},
\act{BeginGracefulShutdown}
  & repaint only (viewport may re-layout)
  & any scrollback emission; any ledger append; frontier motion \\
\addlinespace
\act{Admit}
  & repaint only; new $1$-row slot appears at bottom
  & admission that overflows the viewport; slot for a non-FIFO block \\
\addlinespace
\act{ApplyAllocation}
  & repaint at bridged height $B(\alpha_i,\tau_i)$
  & exceeding the joint reservation \eqref{eq:reservation}; shrinking
    $k{>}2\to 1$ without the $2$-row frame \\
\addlinespace
\act{FinalizeActive}, \act{FinalizeQueued}
  & repaint only; all slots collapse to canonical $1$-row previews
  & scrollback emission; any later mutation of $f_i$ \\
\addlinespace
\act{AppendStable}
  & scroll exactly $\len{\ren{\omega}(\seq{w_{c+1}[n]})}$ rows carrying the
    next stable head row; ledger append of that row
  & streaming any block other than the head $c{+}1$; skipping or reordering
    rows; streaming without row pressure \\
\addlinespace
\act{CompleteAppendOnly}
  & bookkeeping only (frontier $+1$)
  & re-emission of already-streamed rows \\
\addlinespace
\act{RetireSuccess}$(b)$
  & one streamed write $T=F^{\srcRet}_{\mathrm{phys}}(b)\cat\Scr'$
    scrolling exactly $\len{F^{\srcRet}_{\mathrm{phys}}(b)}$ rows
    (Cor.~\ref{cor:retire})
  & retiring a non-contiguous or non-finalized range; re-emitting the emitted
    head prefix; retiring without $\Req$ \\
\addlinespace
\act{RetireFailure}
  & halt; native holds a prefix of the $\srcFail$-tagged batch
  & advancing $c$; mutating the ledger; retrying the write \\
\addlinespace
\act{Resize} ($\Preserve$)
  & record emulator-pushed rows (tagged $\srcRsz$); recanonicalize layout
  & replaying; mutating ledger, frontier, modes, wants, finals, or emitted
    counts \\
\act{Resize} ($\AppendPol$)
  & as $\Preserve$, then schedule replay of $C(c)$ + head prefix at new width
  & same \\
\act{Resize} ($\Rebuild$)
  & clear display epoch ($\Native:=\emptyseq$, $\varepsilon{+}1$), schedule
    replay
  & same \\
\addlinespace
\act{PrepareReplay} + \act{ReplaySynchronousSuccess}
  & paint tail $\RRows[\rcut^{*}{+}1..]$ into blank rows bottom-first
    (no scroll), then \emph{one} synchronous buffered write scrolling exactly
    $\rcut^{*}$ rows
  & interleaving any repaint between preparation and the write; multi-frame
    incremental replay; scrolling more or fewer than $\rcut^{*}$ rows \\
\addlinespace
\act{ReplaySynchronousFailure}
  & halt; native holds a prefix of $\RRows[1..\rcut]$
  & continuing after a partial write \\
\addlinespace
\act{GracefulExit}, \act{DetachExit}
  & optionally scroll the single top viewport row; restore terminal state
  & emitting speculative content of uncommitted blocks; exiting gracefully
    with $c<\created$ or during replay \\
\bottomrule
\end{tabular}
\caption{Conformance obligations per transition.}
\label{tab:conformance}
\end{table}

% ===========================================================================
\section{Formal Theorems and Proofs}\label{sec:theorems}
% ===========================================================================

\subsection{Invariant catalogue}\label{sec:invariants}

\begin{definition}[State invariants]\label{def:inv}
Define the following predicates on states.
\begin{itemize}
\item \textbf{(T)} \emph{Typing}: every component lies in its range of
  Table~\ref{tab:state}.
\item \textbf{(LS)} \emph{Lifecycle shape}: $c\le\created$;
  and for every $i\in\Blocks$,
  \[
    \varphi_i\in\begin{cases}
      \{\Committed\} & i\le c\\
      \{\Queued,\Active,\Finalized\} & c<i\le\created\\
      \{\Absent\} & i>\created,
    \end{cases}
    \qquad
    \mu_i\ne\Undecl \iff i\le\created .
  \]
\item \textbf{(SD)} \emph{Snapshot discipline}: if
  $\varphi_i\in\{\Finalized,\Committed\}$ then $f_i\in\Snap$ and $f_i=w_i$;
  otherwise $f_i=\NoFinal$.
\item \textbf{(ED)} \emph{Emission discipline}: for all $i$,
  $e_i\le\len{w_i}$; $\mu_i\ne\AppendOnly \Rightarrow e_i=0$; and
  $e_i>0 \Rightarrow i=c{+}1 \wedge \varphi_i\in\{\Active,\Finalized\}$.
\item \textbf{(CAP)} \emph{Capacity}: $\AOK(\alpha,\tau)$ at
  $(\varphi,f,e,\omega,h)$.
\item \textbf{(ECH)} \emph{Exact committed history}: equation \eqref{eq:ech},
  $\Ledger=C(c)\cat A(c)$.
\item \textbf{(RS)} \emph{Replay shape}: if $\rmode=\None$ then
  $\rcur=\rend=\rpart=\rcut=0$ and $\neg\rprep$; if $\rmode\ne\None$ then
  $\rcur=1$, $\rend\in 0..c$, $\rpart\le\Lmax$, and: if $\rprep$ then
  $\rcut=\rcut^{*}$ and
  $\len{\RRows[\rcut{+}1..\len{\RRows}]}\le\Room$, else $\rcut=0$.
\item \textbf{(NSS)} \emph{Native source safety}: for every native row
  $(q,o,x,\upsilon)$ in $\Native$:
  \begin{itemize}
  \item $q=\srcRet \Rightarrow o\in[c] \wedge \varphi_o=\Committed$;
  \item $q\in\{\srcApp,\srcRep\} \Rightarrow o\in\Blocks \wedge
    \bigl(o\in[c] \vee (o=c{+}1 \wedge \mu_o=\AppendOnly)\bigr)$;
  \item $q=\srcFail \Rightarrow \stopr=\stWF$;
  \item $q=\srcExit \Rightarrow \neg\run$.
  \end{itemize}
\end{itemize}
Let $J \equiv \mathrm{T}\wedge\mathrm{LS}\wedge\mathrm{SD}\wedge\mathrm{ED}
\wedge\mathrm{CAP}\wedge\mathrm{ECH}\wedge\mathrm{RS}\wedge\mathrm{NSS}$.
\end{definition}

Two further invariants---\emph{no premature history} (NPH) and \emph{screen
capacity} (SC)---are checked independently by TLC but are derivable from $J$;
we prove them as Corollaries~\ref{cor:nph} and~\ref{cor:sc}.

\subsection{Exact capacity and overflow bounds}

\begin{theorem}[Exact capacity and overflow forcing]\label{thm:capacity}
Let a state satisfy \textup{CAP} and \textup{SD}. Then:
\begin{enumerate}
\item[(i)] \textbf{Exactness.} $\Scr$ is well defined, $\len{\Scr}=h$, block
  $i$ owns exactly $\alpha_i$ cells of $\Scr$, exactly $\sigma$ cells equal
  $\Ocell$, and exactly $h-\sum_i\alpha_i-\sigma\ge 0$ cells equal $\Bcell$.
\item[(ii)] \textbf{Zero-height degradation.} If $h=0$ then $\alpha_i=\tau_i=0$
  for all $i$, $\sigma=0$, and $\Scr=\emptyseq$; every transition remains
  well defined.
\item[(iii)] \textbf{Overflow forcing.} If $\ovf$ and $h>0$, then the visible
  blocks are exactly the $h-1$ presented blocks of largest index, and
  necessarily $\alpha_i=\tau_i=1$ for each of them, $\sigma=1$, and $\Scr$
  contains no blank cell: the layout is \emph{forced} to
  $\Scr=\seq{\Ocell}\cat\pc(i_{h-1})\cat\cdots\cat\pc(i_{1})$ where
  $i_{1}>\cdots>i_{h-1}$ enumerate the visible blocks in decreasing order.
\end{enumerate}
\end{theorem}

\begin{proof}
\emph{(i).} Let $A_{\Sigma}=\sum_i\alpha_i$. Since
$\alpha_i\le\max(\alpha_i,\tau_i)$ pointwise, \eqref{eq:reservation} gives
$A_{\Sigma}\le\Res(\alpha,\tau)\le h-\sigma$, so the blank exponent
$h-A_{\Sigma}-\sigma$ in \eqref{eq:screen} is a natural number. The three
factors of \eqref{eq:screen} contribute $h-A_{\Sigma}-\sigma$, $\sigma$, and
$\sum_i\alpha_i$ cells respectively, totalling $h$. Ownership: the slot factor
$\pc(i)^{\alpha_i}$ contributes exactly $\alpha_i$ cells with owner $i$, and
blank/overflow cells carry owner $0\notin\Blocks$. Well-definedness of
$\pc(i)$ for $\alpha_i>0$: by CAP clause~2, $\alpha_i>0$ implies $\vis(i)$,
hence $\pres(i)$, hence $\varphi_i\in\{\Active,\Finalized\}$; if
$\varphi_i=\Active$ then $\ps(i)=\Un_i(w_i)$ with $w_i\in\Snap$; if
$\varphi_i=\Finalized$ then by SD $f_i\in\Snap$, so $\ps(i)=\Un_i(f_i)$ is a
genuine row sequence. The placeholder branch covers an empty rendering.

\emph{(ii).} Let $h=0$. If $\pi=0$ no block is presented, so no block is
visible and CAP clause~2 zeroes all $\alpha,\tau$. If $\pi>0$ then
$\pi>0=h$, i.e.\ $\ovf$; visibility then requires $h>0$, which fails, so again
no block is visible and all allocations are zero; moreover $\sigma=0$ by its
$h>0$ guard. Hence $\Scr=\Bcell^{0}\cat\emptyseq\cat\emptyseq=\emptyseq$.
Every transition guard and effect refers to $\Scr$, $\Room$, and allocation
predicates that are all well defined at $h=0$ (in particular
\act{GracefulExit}/\act{DetachExit} require $k=0$ when $h=0$, and
$\Room=0$ makes any replay scroll its entire frame, $\rcut^{*}=\len{\RRows}$).

\emph{(iii).} Assume $\ovf$ and $h>0$, so $\pi>h\ge 1$ and $\sigma=1$. List
the presented blocks in decreasing index order $i_1>i_2>\cdots>i_\pi$; then
$\newer(i_k)=k-1$, so $\vis(i_k)\iff k-1<h-1 \iff k\le h-1$. Since
$\pi\ge h+1>h-1$, exactly $h-1$ blocks are visible, namely
$i_1,\dots,i_{h-1}$. By CAP clause~1 each visible block has
$\max(\alpha_i,\tau_i)\ge\alpha_i\ge 1$; summing,
$\Res(\alpha,\tau)\ge h-1$, while \eqref{eq:reservation} gives
$\Res(\alpha,\tau)\le h-\sigma=h-1$. Hence every visible block has
$\max(\alpha_i,\tau_i)=1$, so $\alpha_i=1$ and $\tau_i\le 1$; for active
visible blocks CAP requires $\tau_i\ge 1$, and for finalized visible blocks
$\tau_i=\alpha_i=1$. Thus $\alpha_i=\tau_i=1$ on the visible set and $0$
elsewhere; $A_{\Sigma}=h-1$; the blank exponent is $h-(h-1)-1=0$; and
\eqref{eq:screen} degenerates to the stated forced layout, the slots appearing
in ascending block order, i.e.\ the reversed enumeration.
\end{proof}

\begin{corollary}[Canonical admissibility]\label{cor:canon}
For every ambient tuple $(\varphi,f,e,\omega,h)$ (with \textup{SD} for
well-definedness), $\AOK(\Canon,\Canon)$ holds.
\end{corollary}

\begin{proof}
Clauses 1--2 of Definition~\ref{def:aok} hold by construction of $\Canon$,
provided $1\le H$ whenever some block is visible; visibility of any block
implies $h\ge 1$ (if $\neg\ovf$ then $1\le\pi\le h$; if $\ovf$ then
visibility itself demands $h>0$), and $h\le H$. For
\eqref{eq:reservation}: $\Res(\Canon,\Canon)=|\{i:\vis(i)\}|$. If $h=0$ this
is $0=h-\sigma$ by Theorem~\ref{thm:capacity}(ii). If $\ovf$ and $h>0$ it is
$h-1=h-\sigma$ by (iii). If $\neg\ovf$ every presented block is visible, so
the count is $\pi\le h=h-\sigma$.
\end{proof}

Corollary~\ref{cor:canon} is used ubiquitously below: every transition that
recanonicalizes ($\alpha':=\tau':=\Canon$) re-establishes CAP unconditionally,
at whatever successor tuple it evaluates $\Canon$.

\subsection{Reservation dominance and bridge termination}

\begin{lemma}[Reservation dominance; bridge termination]\label{lem:bridge}
Let $a\ne t$ with $a,t\in 0..H$ and let $B$ be the bridge function of
Definition~\ref{def:bridge}. Then:
\begin{enumerate}
\item[(i)] \emph{Dominance:} $\max(B(a,t),t)\le\max(a,t)$; consequently an
  \act{ApplyAllocation} step never increases any summand of
  $\Res(\alpha,\tau)$, and the reservation invariant \eqref{eq:reservation}
  is preserved verbatim (the ambient tuple, hence $\sigma$ and $h$, is
  unchanged by the step).
\item[(ii)] \emph{Positivity:} if $t\ge 1$ then $B(a,t)\ge 1$; hence an active
  visible block never transits through a $0$-height frame.
\item[(iii)] \emph{Termination:} for fixed $t$, the iteration
  $a_0=a$, $a_{n+1}=B(a_n,t)$ reaches $t$ in at most two steps, and each step
  strictly changes the allocation. Precisely: if $a<t$ or $t\ge 2$ or
  $a\le 2$, one step suffices; if $a>2$ and $t=1$, the trajectory is
  $a\to 2\to 1$.
\end{enumerate}
\end{lemma}

\begin{proof}
\emph{(i).} Case $a<t$ (growth): $B=t$ and $\max(t,t)=t=\max(a,t)$. Case
$a>t$: if $a>2\wedge t=1$ then $B=2$ and $\max(2,1)=2<a=\max(a,t)$ since
$a\ge 3$; otherwise $B=t$ and $\max(t,t)=t<a=\max(a,t)$. In all cases the
per-block reservation summand does not increase; all other summands are
untouched, and $\sigma,h$ depend only on $(\varphi,f,e,\omega,h)$, which
\act{ApplyAllocation} leaves unchanged.

\emph{(ii).} If $a<t$ then $B=t\ge 1$; otherwise $B\in\{t,2\}$ and both are
$\ge 1$ when $t\ge 1$.

\emph{(iii).} If $a<t$: $B(a,t)=t$, one step. If $a>t$ and
$\neg(a>2\wedge t=1)$: $B(a,t)=t$, one step. If $a>2$ and $t=1$:
$B(a,1)=2\ne a$ (as $a\ge 3$), then $B(2,1)$ falls in the ``otherwise'' branch
(since $2\not>2$) giving $1$. Each step changes $\alpha_i$: the guard
$\alpha_i\ne\tau_i$ excludes stuttering, and $B(a,t)\ne a$ in every case above.
\end{proof}

\begin{remark}
Dominance is exactly why \act{RequestAllocation} charges
$\max(\alpha_i,\tau^{\dagger}_i)$ at request time: once a target is admitted,
\emph{every} frame of the bridge animation---including the stabilizing
$2$-row frame of a deep shrink---is guaranteed to fit without re-negotiation,
regardless of how apply steps of distinct blocks interleave.
\end{remark}

\subsection{Stable-head emission and ledger monotonicity}

\begin{lemma}[Stable-head emission]\label{lem:head}
Assume $J$. Then:
\begin{enumerate}
\item[(i)] \emph{Prefix immutability.} For every transition and every
  $i\in\Blocks$, if $e_i>0$ then $w_i'[1..e_i]=w_i[1..e_i]$ whenever
  $e_i'\ge e_i$; and the only transitions with $e_i'<e_i$
  (\act{CompleteAppendOnly}, \act{RetireSuccess}) commit block $i$ in the same
  step.
\item[(ii)] \emph{Exactness preservation by streaming.} \act{AppendStable}
  preserves \textup{ECH}: $\Ledger'=C(c)\cat A'(c)$.
\item[(iii)] \emph{No re-emission.} Across any behavior, each position of each
  block's final snapshot contributes at most one ledger row and at most one
  batch of $\srcApp$/$\srcRet$ native rows: rows emitted by
  \act{AppendStable} are excluded from the block's retirement batch
  (Definition~\ref{def:retbatch} slices from $e_{c+1}{+}1$), and after
  commitment $e_i=0$ with block $i$ never again the head.
\item[(iv)] \emph{Ledger monotonicity.} Every transition satisfies
  $\Ledger\pre\Ledger'$.
\end{enumerate}
\end{lemma}

\begin{proof}
\emph{(i).} By ED, $e_i>0$ implies $\mu_i=\AppendOnly$, $i=c{+}1$, and
$\varphi_i\in\{\Active,\Finalized\}$. The transitions that modify $w_i$ are
\act{Update} and the two finalizations, each guarded (for
$\mu_i=\AppendOnly$) by $w_i\pre s$ with $w_i':=s$; a prefix extension
preserves every position $\le e_i\le\len{w_i}$. No other transition writes
$w$. The transitions that decrease $e_i$ are exactly
\act{CompleteAppendOnly} and \act{RetireSuccess}, both of which set
$\varphi_i'=\Committed$ for the affected block(s).

\emph{(ii).} The guard gives $e_{c+1}<\len{w_{c+1}}$ and, writing
$n=e_{c+1}+1$,
\[
  \Ledger' = \Ledger\cat\seq{(c{+}1,w_{c+1}[n])}
           = C(c)\cat A(c)\cat\seq{(c{+}1,w_{c+1}[n])}
           = C(c)\cat\tg_{c+1}\bigl(w_{c+1}[1..n]\bigr),
\]
using ECH for the middle equality and, for the last, that
$A(c)=\tg_{c+1}(w_{c+1}[1..e_{c+1}])$ when $e_{c+1}>0$ and
$A(c)=\emptyseq$ when $e_{c+1}=0$ (in either case the concatenation is the
tagged prefix of length $n$). In the successor, $e_{c+1}'=n>0$, and $\mu$,
$\varphi_{c+1}$, $c$, $\created$ are unchanged, so $\PH'$ holds and
$A'(c)=\tg_{c+1}(w_{c+1}[1..n])$. Hence $\Ledger'=C(c)\cat A'(c)$.

\emph{(iii).} A row of block $i$ enters $\Ledger$ either via
\act{AppendStable} at position $e_i{+}1$ (incrementing $e_i$) or via
\act{RetireSuccess}, whose logical batch for the head starts at position
$e_{c+1}{+}1$ and covers later blocks in full---blocks $j>c{+}1$ have $e_j=0$
by ED. By ECH the ledger is at all times \emph{equal} to
$C(c)\cat A(c)$, in which position $p$ of block $i$ occurs at most once;
monotonicity of $c$ and of $e_{c+1}$ between commits (part (i)) makes the
occurrence permanent. The same slicing governs the physical batches.

\emph{(iv).} Only \act{AppendStable} and \act{RetireSuccess} write $\Ledger$,
both by right-concatenation.
\end{proof}

\subsection{Generic terminal streaming realization}\label{sec:realization}

We now connect the abstract effect ``append $F$ to $\Native$ and repaint the
viewport as $\Scr'$'' to the concrete mechanics of a scrolling terminal.

\begin{definition}[Abstract scrolling terminal]\label{def:terminal}
A \emph{terminal} of height $h\ge 0$ is a pair $(S,V)$ with $S$ a row
sequence (scrollback) and $V$ a row sequence of length $h$ (viewport). The
\emph{print} operation writes a row sequence $T$ starting at viewport row
$r_0\in[h]\cup\{h{+}1\}$ (when $h=0$, $r_0=1$ and every written row scrolls):
rows are written top-to-bottom at consecutive positions; writing at a
position $r\le h$ replaces $V[r]$; when the cursor would advance past row
$h$, the terminal \emph{scrolls}: $S:=S\cat\seq{V[1]}$,
$V:=V[2..h]\cat\seq{\BlRow}$, and the cursor stays at row $h$.
\end{definition}

\begin{lemma}[Streaming realization]\label{lem:stream}
Let $(S,V)$ be a terminal of height $h$, let $F$ be any row sequence and $Q$
any row sequence with $\len{Q}=h$, and let $T=F\cat Q$, $k=\len{F}$. Printing
$T$ from row $r_0=1$ when $k\ge 1$, or overwriting rows $1..h$ when $k=0$,
yields the terminal $(S\cat F,\;Q)$: the scrollback is extended by exactly
$F$, in order, and the final viewport equals $Q$. No row of the pre-state
viewport $V$ enters $S$.
\end{lemma}

\begin{proof}
For $k=0$ the print rewrites $V[1..h]:=Q$ with no cursor overflow; $S$ is
unchanged and $S\cat F=S$.

For $k\ge 1$ we prove by induction on $j\in 0..k$ the invariant: after
writing the first $h+j$ rows of $T$,
\[
  S_j=S\cat T[1..j], \qquad V_j=T[j{+}1\,..\,j{+}h].
\]
\emph{Base $j=0$}: the first $h$ writes fill rows $1..h$ with $T[1..h]$
without overflow, so $S_0=S$ and $V_0=T[1..h]$. \emph{Step}: to write row
$h{+}j{+}1$ of $T$ the cursor advances past row $h$; the terminal scrolls,
moving $V_j[1]=T[j{+}1]$ into scrollback and shifting, then the write places
$T[h{+}j{+}1]$ at row $h$:
\[
  \begin{aligned}
  S_{j+1}&=S_j\cat\seq{T[j{+}1]}=S\cat T[1..j{+}1],\\
  V_{j+1}&=T[j{+}2\,..\,j{+}h]\cat\seq{T[h{+}j{+}1]}
         =T[(j{+}1){+}1\,..\,(j{+}1){+}h].
  \end{aligned}
\]
At $j=k$: $S_k=S\cat T[1..k]=S\cat F$ and
$V_k=T[k{+}1..k{+}h]=Q$. Rows scrolled out are exactly $T[1..k]=F$; every row
of the old viewport is overwritten in the base phase before any scroll, so
none reaches $S$. When $h=0$, $T=F$, every write scrolls, and the claim is
the base-free instance of the same induction.
\end{proof}

\begin{corollary}[Retirement realization]\label{cor:retire}
\act{RetireSuccess}$(b)$ is realized on a terminal displaying $\Scr$ by a
single streamed write of
$T=F^{\srcRet}_{\mathrm{phys}}(b)\cat\Scr'$ (rows taken as the row components
of the tagged sequences), scrolling exactly
$\len{F^{\srcRet}_{\mathrm{phys}}(b)}$ times; likewise \act{AppendStable} with
$F=\ren{\omega}(\seq{w_{c+1}[n]})$ ($1$ or $2$ physical rows), and
\act{GracefulExit}/\act{DetachExit} with $k=1$ and $F=\seq{\Scr[1]}$ followed
by a blanked viewport. In each case the scrollback extension equals, row for
row and in order, the native rows appended by the transition, and no
speculative viewport content leaks into scrollback.
\end{corollary}

\begin{proof}
Instantiate Lemma~\ref{lem:stream} with $Q=\Scr'$, which has length $h$ by
Theorem~\ref{thm:capacity}(i) applied in the successor state (CAP holds there
by Corollary~\ref{cor:canon}). The transition's native append is by
definition the tagged image of $F$. The final claim is the lemma's ``no row of
$V$ enters $S$''.
\end{proof}

\begin{remark}[Fail-stop realization]
\act{RetireFailure}$(b,k)$ is the same write interrupted after $k$ physical
rows: by Lemma~\ref{lem:stream}'s induction, after $\min(k,\dots)$ writes the
scrollback extension is exactly a \emph{prefix} of
$F^{\srcFail}_{\mathrm{phys}}(b)$---never a row of $\Scr'$, never out of
order. Because the abstract transition does not advance $c$ or touch
$\Ledger$, and $\run'=\mathrm{false}$ disables every subsequent transition
(Theorem~\ref{thm:actions}(vii)), the physical scrollback ends with a clean,
attributable, immutable prefix of finalized content.
\end{remark}

\subsection{Bottom-first synchronous replay}\label{sec:replayproof}

\begin{lemma}[Bottom-first synchronous replay invariant]\label{lem:replay}
Assume $J$ and $\Repl$. Then:
\begin{enumerate}
\item[(i)] \emph{Split correctness.} After \act{PrepareReplay},
  $\rcut=\rcut^{*}$ and the prepared tail satisfies
  $\len{\RRows[\rcut{+}1..\len{\RRows}]}=\min(\len{\RRows},\Room)\le\Room$.
\item[(ii)] \emph{History-neutral tail paint.} Painting the prepared tail
  into the blank region of $\Scr$ (bottom-anchored, by absolute cursor
  addressing) writes only rows at viewport positions occupied by $\Bcell$
  cells and causes no scroll; hence it extends neither $\Native$ nor
  $\Ledger$.
\item[(iii)] \emph{Single-write scroll exactness.} The subsequent synchronous
  buffered write is the streaming realization
  (Lemma~\ref{lem:stream}) of $F=\RRows[1..\rcut]$ with final viewport the
  replay-completed frame; it scrolls exactly $\rcut$ rows, and the native
  extension is exactly $\RRows[1..\rcut]$---as specified by
  \act{ReplaySynchronousSuccess}.
\item[(iv)] \emph{Gate necessity and atomicity.} Between \act{PrepareReplay}
  and the write, the scheduler gate of Definition~\ref{def:spec} admits no
  other transition; this is necessary: absent the gate, an interleaved
  \act{Admit}, \act{Update}, \act{FinalizeQueued}, or allocation step can
  change $\Room$ or $\len{\RRows}$ and falsify $\rcut=\rcut^{*}$, so the
  invariant \textup{RS} is stable only under the gate. Consequently at most
  one frame is ever presented per replay: no multi-frame flicker and no torn
  intermediate state is observable.
\item[(v)] \emph{Failure safety.} \act{ReplaySynchronousFailure}$(k)$ leaves
  scrollback extended by the prefix $\RRows[1..k]$, $k\le\rcut$, of
  re-rendered \emph{committed} rows (and stable head prefix); the ledger and
  frontier are untouched and the host halts.
\end{enumerate}
\end{lemma}

\begin{proof}
\emph{(i).} \act{PrepareReplay} sets $\rcut':=\rcut^{*}$ by definition. Then
\[
  \len{\RRows}-\rcut^{*}
  =\len{\RRows}-\max(0,\len{\RRows}-\Room)
  =\min(\len{\RRows},\Room)\le\Room .
\]

\emph{(ii).} By (i) the tail occupies at most $\Room$ rows and $\Room$ counts
exactly the $\Bcell$ positions of $\Scr$
(Theorem~\ref{thm:capacity}(i) makes this count well defined and the blank
block contiguous at the top of $\Scr$, positions
$1..\,h-A_{\Sigma}-\sigma$). Writing $\le\Room$ rows at absolute positions
within $[1..h]$ never advances the cursor past row $h$, so by
Definition~\ref{def:terminal} no scroll occurs and $S$ is unchanged. The
transition \act{PrepareReplay} indeed changes neither $\Native$ nor
$\Ledger$.

\emph{(iii).} Instantiate Lemma~\ref{lem:stream} with $F=\RRows[1..\rcut]$
and $Q$ the post-replay physical frame (the $h$-row image containing the
painted tail and the live slots). The scrollback extension is exactly
$F$, matching the abstract effect
$\Native':=\Native\cat\RRows[1..\rcut]$ of \act{ReplaySynchronousSuccess};
$\Ledger$, $c$, $w$, $f$, $e$ are unchanged, so replay is logically neutral.

\emph{(iv).} By Definition~\ref{def:spec}, when $\rprep$ holds the next-state
relation is the disjunction of the two replay writes only. For necessity,
observe that $\rcut^{*}$ is a function of $\len{\RRows}$ and $\Room$;
$\Room$ is a function of $\alpha$ and $\sigma$, both of which are changed by
\act{Admit} (new $1$-row slot), \act{ApplyAllocation}, \act{FinalizeQueued}
(possible new presented block, possibly overflow), and \act{Update} on the
head can change $\len{\RRows}$ via $w_{\rend+1}$\,---\,so RS's clause
$\rprep\Rightarrow\rcut=\rcut^{*}$ would not be preserved by those steps.
Physically, the gate is the statement that the implementation assembles the
remainder $F\cat Q$ into one buffer and issues it as one synchronous write:
no repaint can be scheduled between preparation and completion, so the
display transitions from the pre-replay frame directly to the completed
frame.

\emph{(v).} Immediate from the transition's definition and
Lemma~\ref{lem:stream}'s prefix property; NSS's $\srcRep$ clause (proved in
Theorem~\ref{thm:safety}) attributes every such row to committed content or
the append-only head.
\end{proof}

\begin{remark}[Bottom-first order]
The order---tail first, scroll-write second---is what makes the \emph{visible}
viewport correct at every instant: the tail paint (ii) fills blank rows that
carry no information, and the single write (iii) atomically supplies the
scrolled prefix and the final frame. Painting top-first (prefix before tail)
would either require $\rcut$ speculative scrolls before the tail exists on
screen (tearing: transient frames showing history without its continuation)
or force multi-frame incremental replay (flicker). The model's gate plus
split renders both failure modes unrepresentable.
\end{remark}

\subsection{Global inductive safety}\label{sec:safety}

\begin{theorem}[Global inductive safety]\label{thm:safety}
$J$ (Definition~\ref{def:inv}) is an inductive invariant of
$\mathrm{Spec}$: the initial state satisfies $J$, and every transition of
Section~\ref{sec:protocol}, fired in any state satisfying $J$, yields a state
satisfying $J$. Hence $\mathrm{Spec}\models\Box J$.
\end{theorem}

\begin{proof}
\emph{Initiality.} In the state of Definition~\ref{def:init}: T holds by
inspection; LS with $c=\created=0$; SD with all phases $\Absent$ and
$f\equiv\NoFinal$; ED with $e\equiv 0$; CAP because no block is presented, so
$\Canon\equiv 0=\alpha=\tau$ and $\Res+\sigma=0\le H=h$
(Corollary~\ref{cor:canon}); ECH because
$\Ledger=\emptyseq=C(0)\cat A(0)$ ($\PH$ fails at $c=\created$); RS by the
$\None$ clause; NSS vacuously ($\Native=\emptyseq$).

\emph{Preservation.} Fix a state satisfying $J$ and a transition. We argue
each conjunct; components not written by the transition preserve the
corresponding clauses trivially, so we treat only the writers.

\textbf{T.} All effects assign values within the declared ranges:
phases/modes from their alphabets; $e_{c+1}'=e_{c+1}+1\le\len{w_{c+1}}\le
\Lmax$ in \act{AppendStable} (guard); $\alpha',\tau'$ from guards or
$\Canon\in\{0,1\}^{\Blocks}$; $c'=b\le N$; $\nu'\le\MaxRes$ (guard) and
$\varepsilon'\le\nu'\le\MaxRes$ (an epoch increment occurs only within a
resize); $\rcut'=\rcut^{*}\le\len{\RRows}\le
2\bigl(c\,\Lmax+\Lmax\bigr)\le 2N\Lmax=\Fmax$ using
Proposition~\ref{prop:render}(i); ledger/native rows are constructed by
$\tg$/$\ntg$ over their row types.

\textbf{LS.} \act{Create} sets $\varphi_{\created+1}$ from $\Absent$ to
$\Queued$ with a mode, extending the created range contiguously by one; the
bands shift accordingly. \act{Admit}, \act{Update}, the finalizations, and
\act{BeginGracefulShutdown} keep every touched block inside the band
prescribed for its index: indices $\le c$ are untouched or (shutdown) set to
$\Committed$, indices in $(c..\created]$ move within
$\{\Queued,\Active,\Finalized\}$. \act{CompleteAppendOnly} and
\act{RetireSuccess}$(b)$ set $\varphi_j'=\Committed$ exactly for
$j\le c'\in\{c{+}1,b\}$ and leave later blocks in place; $c'\le\created$
because the guard requires those blocks finalized, hence created. No
transition unsets a mode or reverts a phase.

\textbf{SD.} Writers of $f$: the finalizations set $f_i'=s=w_i'$ with
$\varphi_i'=\Finalized$ (and by LS the block was not finalized/committed
before, so no frozen final is overwritten); \act{BeginGracefulShutdown} sets
$f_i'=w_i$ exactly for blocks moving $\Queued/\Active\to\Finalized$
(for which $w$ is unchanged, giving $f_i'=w_i'$) and preserves $f_i$ for
already finalized/committed blocks, where SD gave $f_i=w_i$ and $w$ is
unchanged. Writers of $w$ with a frozen final: none---\act{Update} requires
$\varphi_i\in\{\Queued,\Active\}$ (so $f_i=\NoFinal$), and the finalizations
set both together. Commitment (\act{CompleteAppendOnly},
\act{RetireSuccess}) changes phases $\Finalized\to\Committed$ with $f,w$
untouched.

\textbf{ED.} Writers of $e$: \act{AppendStable} increments $e_{c+1}$ within
$\len{w_{c+1}}$ (guard), for a block with $\mu=\AppendOnly$,
$\varphi\in\{\Active,\Finalized\}$, index $c{+}1$, and leaves $c$ fixed;
\act{CompleteAppendOnly} and \act{RetireSuccess} zero $e_j$ for all $j\le c'$
while advancing $c$---by ED any $i$ with $e_i>0$ satisfied $i=c{+}1\le c'$,
so no positive $e$ survives at an index $\ne c'{+}1$; indices $>c'$ retain
$e=0$. Writers of $w$: for $\mu_i=\AppendOnly$ the prefix guards give
$\len{w_i'}\ge\len{w_i}\ge e_i$; for $\mu_i=\Mutable$, ED gives $e_i=0$.
Writers of $\varphi$ at an index with $e_i>0$: only
\act{FinalizeActive} ($\Active\to\Finalized$, allowed) and
\act{BeginGracefulShutdown} (by ED the block with $e_i>0$ has
$\varphi_i\in\{\Active,\Finalized\}$ and $i=c{+}1>c$, so it moves to
$\Finalized$, allowed).

\textbf{CAP.} \act{Admit}, \act{RequestAllocation}, and
\act{ApplyAllocation} carry $\AOK$ at the successor configuration explicitly
in their guards (for \act{ApplyAllocation} the guard is moreover always
satisfiable along a bridge by Lemma~\ref{lem:bridge}(i)). \act{Update}
changes only $w$, on which $\AOK$ does not depend (visibility uses
$\varphi,f,e,\omega,h$ only). \act{BeginFlush}, \act{PrepareReplay}, the
replay writes, \act{RetireFailure}, and the exits change none of
$\alpha,\tau,\varphi,f,e,\omega,h$. Every remaining writer
(\act{FinalizeActive}, \act{FinalizeQueued}, \act{AppendStable},
\act{CompleteAppendOnly}, \act{RetireSuccess}, \act{Resize},
\act{BeginGracefulShutdown}) sets $\alpha'=\tau'=\Canon$ evaluated at its
successor tuple, and Corollary~\ref{cor:canon} yields $\AOK$ there.

\textbf{ECH.} Writers of any of $\Ledger,c,f,w,e,\varphi,\mu$:
\begin{itemize}
\item \act{Create}: $\created'=\created+1$. If $c<\created$, the head block
  $c{+}1$ is untouched, so $\PH$ and $A(c)$ are unchanged. If $c=\created$,
  then $\PH$ was false; afterwards $e_{c+1}=0$ (ED), so $\PH'$ is false.
  Either way $\Ledger'=\Ledger=C(c)\cat A(c)$.
\item \act{Admit}, \act{ApplyAllocation}, \act{RequestAllocation},

  \act{BeginFlush}, \act{Resize}, \act{PrepareReplay}, the replay writes,
  failures, exits: none of $\Ledger,c,f,w,e$ change, and $\PH$ depends only
  on $(c,\created,\mu_{c+1},\varphi_{c+1},e_{c+1})$; \act{Admit} can change
  $\varphi_{c+1}$ only $\Queued\to\Active$, which affects $\PH$ only when
  $e_{c+1}>0$, impossible for a queued block by ED.
\item \act{Update}$(i,s)$: only $w_i$ changes. If $i\ne c{+}1$ or $\PH$
  fails with $e_{c+1}=0$, $A(c)$ is unaffected. If $i=c{+}1$ and
  $e_{c+1}>0$, ED forces $\mu_i=\AppendOnly$, the guard gives $w_i\pre s$,
  hence $s[1..e_i]=w_i[1..e_i]$ and $A'(c)=A(c)$.
\item Finalizations: $f_i$ is set for $i>c$, so $C(c)$ is unchanged; $w_i$
  extends by prefix in the append-only case (as above) and $e_i=0$ in the
  mutable case; $\varphi_i'=\Finalized$ keeps
  $\PH$'s phase disjunct satisfied when $i=c{+}1$.
\item \act{AppendStable}: Lemma~\ref{lem:head}(ii).
\item \act{CompleteAppendOnly}: the guard gives $e_{c+1}=\len{f_{c+1}}$ and
  SD gives $f_{c+1}=w_{c+1}$, so
  $A(c)=\tg_{c+1}(w_{c+1}[1..e_{c+1}])=\tg_{c+1}(f_{c+1})$ (also when
  $\len{f_{c+1}}=0$, both sides empty). Hence
  $\Ledger=C(c)\cat\tg_{c+1}(f_{c+1})=C(c{+}1)=C(c')$. In the successor,
  every $e'_j=0$ for $j\le c'$ and ED excludes $e_{c'+1}>0$ (only index
  $c{+}1$ could be positive), so $\PH'$ fails and $A'(c')=\emptyseq$.
\item \act{RetireSuccess}$(b)$: all blocks $c{+}1..b$ are finalized, so
  SD gives $f_j=w_j$ on the range. If $e_{c+1}>0$ then
  $A(c)\cat\tg_{c+1}(f_{c+1}[e_{c+1}{+}1..])=\tg_{c+1}(f_{c+1})$; if
  $e_{c+1}=0$ then $A(c)=\emptyseq$ and the head batch is all of
  $\tg_{c+1}(f_{c+1})$. In both cases
  \[
    \Ledger'=C(c)\cat A(c)\cat F_{\log}(b)
    = C(c)\cat\Cat_{j=c+1}^{b}\tg_j(f_j)=C(b)=C(c').
  \]
  As above, all $e'_j$ vanish for $j\le b$ and $e_{b+1}=0$ by ED, so
  $A'(c')=\emptyseq$.
\item \act{BeginGracefulShutdown}: $\Ledger,c,w,e$ unchanged; $f_j$ for
  $j\le c$ unchanged, so $C(c)$ unchanged; if $\PH$ held, the head's phase
  moves $\Active\mapsto\Finalized$ or stays $\Finalized$, so $\PH'$ holds
  and $A$ is unchanged ($w,e$ fixed); if $\PH$ failed, it fails after
  (a positive $e_{c+1}$ cannot appear).
\end{itemize}

\textbf{RS.} \act{Resize}: in the $\mathit{begin}$ case it establishes the
$\Repl$ clause directly ($\rcur'=1$, $\rend'=c\in 0..c$,
$\rpart'\le\Lmax$ by T, $\rprep'$ false, $\rcut'=0$); continuing an existing
replay it preserves the window (still $\rend\le c$ since $c$ is unchanged)
and resets $\rprep,\rcut$ to the unprepared shape; otherwise it writes the
$\None$ shape. \act{PrepareReplay} establishes the prepared clause:
$\rcut'=\rcut^{*}$ and the tail bound by Lemma~\ref{lem:replay}(i).
\emph{Stability of the prepared clause}: by the scheduler gate
(Definition~\ref{def:spec}) the only transitions from a prepared state are
the two replay writes; success restores the $\None$ shape and failure
freezes the system, so no step can invalidate $\rcut=\rcut^{*}$
(Lemma~\ref{lem:replay}(iv)). While unprepared and replaying, transitions
other than \act{Resize}/\act{PrepareReplay} leave the replay components
untouched, and all $c$-advancing transitions
(\act{CompleteAppendOnly}, \act{RetireSuccess}) are disabled by their
$\neg\Repl$ guards, preserving $\rend\le c$.

\textbf{NSS.} New native rows, by transition: \act{AppendStable} appends
$\srcApp$ rows with owner $c{+}1$ and $\mu_{c+1}=\AppendOnly$ (guard);
\act{RetireSuccess}$(b)$ appends $\srcRet$ rows with owners in
$c{+}1..b\subseteq[c']$, all committed in the successor;
\act{Resize} appends $\srcRsz$ rows (no NSS obligation) or clears $\Native$
(vacuous); the replay writes append $\srcRep$ rows whose owners lie in
$[\rend]\cup\{\rend{+}1\}$ with $\rend\le c$ and, when $\rpart>0$, owner
$\rend{+}1=c{+}1$ with $\mu_{c+1}=\AppendOnly$---the latter because
$\rpart>0$ was set at a resize at which $\PH$ held, $\mu$ is immutable after
creation, and $c$ cannot advance while $\Repl$ holds, so the identity
$\rend=c$ persists from that resize; \act{RetireFailure} and
\act{ReplaySynchronousFailure} append $\srcFail$- and $\srcRep$-tagged rows
respectively while setting $\stopr'=\stWF$ and $\run'=\mathrm{false}$; the
exits append $\srcExit$ rows while setting $\run'=\mathrm{false}$.
Old rows: the clauses for $\srcRet$, $\srcApp$, $\srcRep$ mention $c$,
$\varphi$, $\mu$, all of which evolve monotonically in the required
direction---$c$ never decreases, committed blocks stay committed, modes are
immutable, and an owner $o=c{+}1$ of an $\srcApp$/$\srcRep$ row moves into
$[c']$ whenever the frontier passes it. The $\srcFail$/$\srcExit$ clauses
are preserved because $\stopr$ and $\run$ change only in halting transitions,
after which (Theorem~\ref{thm:actions}(vii)) the state is frozen.
\end{proof}

\begin{corollary}[No premature native or ledger rows]\label{cor:nph}
In every reachable state, every ledger row is owned either by a committed
block $o\in[c]$ (with $\varphi_o=\Committed$) or, when $\PH$ holds, by the
head $c{+}1$; and every $\srcApp$/$\srcRet$/$\srcRep$ native row is owned by
a committed block or by the append-only head. In particular no row of a
queued, merely-active mutable, or unfinalized non-head block ever reaches
either history.
\end{corollary}

\begin{proof}
By ECH, $\Ledger=C(c)\cat A(c)$: rows of $C(c)$ have owners $j\in[c]$, which
LS makes committed; rows of $A(c)$ exist only under $\PH$ and carry owner
$c{+}1$. The native claim is NSS. Both hold in all reachable states by
Theorem~\ref{thm:safety}.
\end{proof}

\begin{corollary}[Screen capacity]\label{cor:sc}
In every reachable state, $\Scr$ is an $h$-row cell sequence in which block
$i$ owns exactly $\alpha_i$ rows, exactly $\sigma$ rows carry the overflow
marker, and exactly $h-\sum_i\alpha_i-\sigma$ rows are blank.
\end{corollary}

\begin{proof}
Theorem~\ref{thm:capacity}(i) applied under $\Box J$.
\end{proof}

\begin{theorem}[Action safety]\label{thm:actions}
Every step of $\mathrm{Spec}$ (including stuttering) satisfies:
\begin{enumerate}
\item[(i)] \emph{Ledger monotonicity:} $\Ledger\pre\Ledger'$.
\item[(ii)] \emph{Epoch-scoped native monotonicity:} either
  $\varepsilon'=\varepsilon$ and $\Native\pre\Native'$, or
  $\varepsilon'=\varepsilon+1$ and $\Native'=\emptyseq$.
\item[(iii)] \emph{Final immutability:} if
  $\varphi_i\in\{\Finalized,\Committed\}$ then $f_i'=f_i$.
\item[(iv)] \emph{Append-only monotonicity:} if $\mu_i=\AppendOnly$ and
  $\varphi_i\in\{\Queued,\Active\}$ then $w_i\pre w_i'$.
\item[(v)] \emph{Resize logical-neutrality:} if $\omega'\ne\omega$ or
  $h'\ne h$ then $\Ledger'=\Ledger$, $c'=c$, $\mu'=\mu$, $w'=w$, $f'=f$,
  $e'=e$.
\item[(vi)] \emph{Fail-stop:} $\Box(\stopr=\stWF \Rightarrow \neg\run)$.
\item[(vii)] \emph{Stopped quiescence:} if $\neg\run$ then the step is a
  stutter (all components unchanged).
\end{enumerate}
\end{theorem}

\begin{proof}
(i) Lemma~\ref{lem:head}(iv); stutters are trivial.
(ii) The writers of $\Native$ are \act{AppendStable}, \act{RetireSuccess},
\act{RetireFailure}, \act{Resize}, the replay writes, and the exits; all
append except \act{Resize} with $p^{\dagger}=\Rebuild$, which simultaneously
sets $\Native'=\emptyseq$ and $\varepsilon'=\varepsilon+1$; no other
transition writes $\varepsilon$.
(iii) The finalizations write $f_i$ only for $\varphi_i\in\{\Active,\Queued\}$;
\act{BeginGracefulShutdown} preserves $f_i$ whenever
$i\le c$ or $\varphi_i=\Finalized$, which by LS covers all finalized and
committed blocks.
(iv) The writers of $w_i$ are \act{Update} and the finalizations; for
$\mu_i=\AppendOnly$ each is guarded by $w_i\pre s$.
(v) Only \act{Resize} writes $\omega$ or $h$, and its effect leaves the six
listed components untouched.
(vi) $\stopr$ is set to $\stWF$ only by \act{RetireFailure} and
\act{ReplaySynchronousFailure}, both of which set $\run'=\mathrm{false}$ in
the same step; by (vii) no later step changes either component; initially
$\stopr=\stRun$.
(vii) Every transition's guard conjoins $\run$; hence from a state with
$\neg\run$, $\mathrm{Next}$ is disabled and $\Box[\mathrm{Next}]_{vars}$
permits only stuttering. (Weak fairness of the five named actions is
vacuous there, since none is enabled.)
\end{proof}

\begin{remark}[Fail-stop precludes duplication and reordering]
Combining Theorem~\ref{thm:actions}(vi,vii) with
Theorem~\ref{thm:safety}: after a partial write failure, $c$ and $\Ledger$
are exactly what they were before the failed batch, the native scrollback
ends in a $\srcFail$- or $\srcRep$-tagged \emph{prefix} of immutable
finalized/stable rows, and no further transition exists. There is thus no
state from which a retry could re-emit rows, and no interleaving in which a
younger block's rows precede an older block's in either history.
\end{remark}

\subsection{Conditional progress and liveness}\label{sec:liveness}

Recall the weak-fairness conjuncts of Definition~\ref{def:spec} and write
$\leadsto$ for ``leads to''. Define
\[
  \mathrm{AllFin} \iff \forall i\in[\created]:
    \varphi_i\in\{\Finalized,\Committed\},
  \qquad
  \mathrm{AllCom} \iff c=\created \wedge \Ledger=C(c).
\]

\begin{lemma}[Resize exhaustion]\label{lem:resize-exhaust}
In every behavior of $\mathrm{Spec}$, only finitely many \act{Resize} steps
occur; after the last one, $\omega$ and $h$ are constant and no new replay
begins.
\end{lemma}
\begin{proof}
$\nu$ increases by one at each \act{Resize} and at no other step, and
$\nu\le\MaxRes$ by T. \act{Resize} is the only transition that sets
$\rmode$ to a non-$\None$ value.
\end{proof}

\begin{theorem}[Replay termination]\label{thm:replay-live}
$\mathrm{Spec}\models
(\Repl \wedge \run) \leadsto (\neg\Repl \vee \neg\run)$.
\end{theorem}

\begin{proof}
Assume, toward contradiction, a behavior with a suffix on which
$\Repl\wedge\run$ holds forever. By Lemma~\ref{lem:resize-exhaust} pass to a
further suffix with no \act{Resize} steps. On this suffix $\rmode$ can only
be changed by \act{ReplaySynchronousSuccess} (to $\None$, contradiction) and
$\run$ only by a halting transition (contradiction), so suppose neither ever
fires. Two cases on $\rprep$:

If $\rprep$ is false at some state of the suffix, then \act{PrepareReplay} is
enabled there ($\run\wedge\Repl\wedge\neg\rprep$) and remains enabled at
every subsequent state until taken: the only transitions that modify
$\rprep$ or $\rmode$ are \act{Resize} (excluded), \act{PrepareReplay}
itself, and the replay writes (which require $\rprep$). By
$\mathrm{WF}(\act{PrepareReplay})$ it is eventually taken, after which
$\rprep$ holds.

Once $\rprep$ holds, the scheduler gate restricts $\mathrm{Next}$ to the two
replay writes; the action \act{ReplaySynchronousSuccess} is then enabled
($\run\wedge\Repl\wedge\rprep$) and stays enabled while neither write fires
(no other transition can execute). By
$\mathrm{WF}(\act{ReplaySynchronousSuccess})$, eventually the success fires
(making $\rmode=\None$) or the failure fires first (making
$\neg\run$)---either contradicts the assumption.
\end{proof}

\begin{theorem}[Flush-driven completion]\label{thm:flush-live}
$\mathrm{Spec}\models
(\mathrm{AllFin}\wedge\flush\wedge\shut\wedge\run\wedge\neg\Repl)
\leadsto (\mathrm{AllCom}\vee\neg\run)$.
\end{theorem}

\begin{proof}
Let $s_0$ satisfy the antecedent and consider the behavior from $s_0$,
assuming $\run$ persists (otherwise the consequent holds). Under $\shut$
(which no transition unsets), the transitions \act{Create}, \act{Admit},
\act{Update}, \act{RequestAllocation}, \act{ApplyAllocation}, both
finalizations, \act{AppendStable}, \act{Resize},
\act{BeginGracefulShutdown}, and \act{DetachExit} are disabled by their
$\neg\shut$ guards. In particular no new replay can begin
(only \act{Resize} begins one), so $\neg\Repl$ persists; $\created$ is
frozen; $\mathrm{AllFin}$ persists (phases only move toward $\Committed$);
and $\flush$ persists, so $\Req$ holds at every state.

Claim: while $c<\created$, the action
$\act{RetireSuccess}^{\exists}$ is enabled. Indeed $b=\created$ satisfies
$b\in(c{+}1)..N$, the range $c{+}1..b$ is finalized
($\mathrm{AllFin}$, minus committed prefix), $\Req$ holds, and
$\neg\Repl\wedge\run$ hold. The only enabled non-stuttering transitions from
such states are \act{RetireSuccess}, \act{RetireFailure},
\act{CompleteAppendOnly}, \act{BeginFlush} (disabled: $\flush$ already), and
\act{GracefulExit} (disabled while $c<\created$). If \act{RetireFailure}
fires, $\neg\run$ and we are done. \act{CompleteAppendOnly} strictly
increases $c$. Since $\act{RetireSuccess}^{\exists}$ remains enabled at every
state with $c<\created$, weak fairness guarantees that, unless $c$ reaches
$\created$ by \act{CompleteAppendOnly} steps first, a
\act{RetireSuccess}$(b)$ step occurs, which sets $c'=b\ge c+1$. Each
frontier-advancing step increases $c$ by at least one and $c\le N$, so after
at most $N$ such steps $c=\created$. At that point $\PH$ fails
(its conjunct $c<\created$), so by ECH
$\Ledger=C(c)\cat\emptyseq=C(c)$: exactly $\mathrm{AllCom}$.
\end{proof}

\begin{theorem}[Pressure-driven retirement]\label{thm:pressure-live}
Let $i\in\Blocks$. Suppose a reachable state $s$ satisfies
\[
  \run \wedge \neg\Repl \wedge c=i{-}1 \wedge \varphi_i=\Finalized
  \wedge \bigl(\flush \vee \created-c\ge\MaxLive\bigr).
\]
Then $\mathrm{Spec}\models$ from $s$, eventually $c\ge i \vee \neg\run$.
\end{theorem}

\begin{proof}
Assume $\run$ persists and $c=i{-}1$ persists (otherwise, since $c$ is
monotone, $c\ge i$ eventually and we are done). Then: $\flush$, once true,
is stable; and $\created-c$ is nondecreasing while $c$ is fixed
($\created$ is nondecreasing), so the disjunct that held at $s$ holds
forever, hence $\Req$ holds forever. $\varphi_i=\Finalized$ persists while
$c<i$ (a finalized block changes phase only by committing, which forces
$c\ge i$). By Lemma~\ref{lem:resize-exhaust} and
Theorem~\ref{thm:replay-live}, eventually a suffix begins on which
$\neg\Repl$ holds forever. On that suffix
$\act{RetireSuccess}^{\exists}$ is continuously enabled (witness $b=i$:
$\mathrm{FinalizedRange}$ is the single block $i$, $\Req$ holds,
$\neg\Repl\wedge\run$ hold), so by weak fairness some
\act{RetireSuccess}$(b)$ with $b\ge c+1=i$ fires---contradicting the
persistence of $c=i{-}1$. (A \act{RetireFailure} instead yields $\neg\run$.)
\end{proof}

\begin{remark}[Sharpness: row pressure alone does not force retirement]
\label{rem:sharp}
Theorem~\ref{thm:pressure-live} conditions on the \emph{persistent} pressure
disjuncts ($\flush$, or count pressure $\created-c\ge\MaxLive$, which is
monotone while the frontier is fixed). Pure \emph{row} pressure $D>h$ is not
persistent and does not suffice in general: with $\MaxLive=4$, $N=3$, a
finalized head ($f_1=\seq{a}$), an active mutable block~$2$, and a queued
block~$3$ under $h=1$, we have $D=3>1$; a subsequent \act{Update} shrinking
$w_2$ and a $\Preserve$ resize to $h'=3$ give $D=3\le 3$, extinguishing
$\Press$ with $\created-c=3<\MaxLive$ and no flush---after which
$\act{RetireSuccess}$ is disabled forever and an infinite stutter-fair
suffix keeps $c=0$. The mechanically checked property
\tla{QueuedPressureRetirement} (\S\ref{sec:tlc}) adds the hypothesis
$\exists j:\varphi_j=\Queued$ and is checked at $\MaxLive=2$; there the
hypothesis \emph{implies} persistent count pressure---a queued block $j$ and
the finalized head $i$ are distinct uncommitted created blocks, so
$\created-c\ge 2=\MaxLive$---and the property is an instance of
Theorem~\ref{thm:pressure-live}. This is the precise sense in which queued
demand guarantees that older finished work is swept out of the viewport.
\end{remark}

% ===========================================================================
\section{Mechanical Verification in TLA\texorpdfstring{\textsuperscript{+}}{+}}
\label{sec:tlc}
% ===========================================================================

The entire development of \S\ref{sec:model}--\S\ref{sec:theorems} is a
faithful mathematical reading of the module \texttt{ElasticSlots.tla}; this
section records the correspondence and the exhaustive model-checking
evidence.

\subsection{Transition correspondence}

Table~\ref{tab:actions} maps each transition to its
TLA\textsuperscript{+} action. The scheduler gate of
Definition~\ref{def:spec} is the \tla{IF replayPrepared} branch of
\tla{Next}; the five weak-fairness conjuncts appear verbatim in \tla{Spec}.
Failure nondeterminism is bounded by
\tla{MaxFailureRows} $=\Fmax=2N\Lmax$.

\begin{table}[htbp]
\centering\footnotesize\setlength{\tabcolsep}{3.5pt}
\begin{tabular}{@{}ll@{\hspace{1em}}ll@{}}
\toprule
This paper & TLA\textsuperscript{+} action &
This paper & TLA\textsuperscript{+} action \\
\midrule
\act{Create}$(m)$ & \tla{Create} &
\act{RetireSuccess}$(b)$ & \tla{RetireSuccess} \\
\act{Admit}$(i)$ & \tla{Admit} &
\act{RetireFailure}$(b,k)$ & \tla{RetireFailure} \\
\act{Update}$(i,s)$ & \tla{Update} &
\act{Resize}$(\omega^{\dagger},h^{\dagger},p,k)$ & \tla{Resize} \\
\act{RequestAllocation}$(\tau^{\dagger})$ & \tla{RequestAllocation}
  & \act{PrepareReplay} & \tla{PrepareReplay} \\
\act{ApplyAllocation}$(i)$ & \tla{ApplyAllocation} &
\act{ReplaySynchronousSuccess} & \tla{ReplaySynchronousSuccess} \\
\act{FinalizeActive}$(i,s)$ & \tla{FinalizeActive} &
\act{ReplaySynchronousFailure}$(k)$ & \tla{ReplaySynchronousFailure} \\
\act{FinalizeQueued}$(i,s)$ & \tla{FinalizeQueued} &
\act{BeginGracefulShutdown} & \tla{BeginGracefulShutdown} \\
\act{AppendStable} & \tla{AppendStable} &
\act{GracefulExit}$(k)$ & \tla{GracefulExit} \\
\act{CompleteAppendOnly} & \tla{CompleteAppendOnly} &
\act{DetachExit}$(k)$ & \tla{DetachExit} \\
\act{BeginFlush} & \tla{BeginFlush} & & \\
\bottomrule
\end{tabular}
\caption{Transition-to-action correspondence (TLA\textsuperscript{+}
parameter lists---\tla{declaration}, \tla{snapshot}, \tla{newTarget},
\tla{batchEnd}, \tla{count}, \tla{newWidth}, \tla{newHeight},
\tla{resizePolicy}, \tla{pushed}, \tla{push}---mirror the mathematical
parameters positionally).}
\label{tab:actions}
\end{table}

Likewise for the derived operators: $\ren{\omega}$ is
\tla{Render}/\tla{DoubleRows}; $\tg$/$\ntg$ are \tla{Tag}/\tla{NativeTag};
$C(k)$ is \tla{CommittedRows}; $A(c)$ is \tla{PartialHeadRows} (with $\PH$
as \tla{PartialHeadExists}); $F_{\log},F^{q}_{\mathrm{phys}}$ are
\tla{RetirementRows}/\tla{NativeRetirementRows}; $\pres,\vis,\sigma$ are
\tla{Presented}, \tla{VisiblePresented}, \tla{SummaryRows}; $\AOK$ is
\tla{AllocationStateOK}; $\Canon$ is \tla{CanonicalAllocation}; $B$ is
\tla{BridgeHeight}; $\dm,D$ are \tla{SnapshotHeight}/\tla{FullRows};
$\Press,\Req$ are \tla{Pressure}/\tla{RetirementRequested}; $\Scr$ is
\tla{Screen}; $\RRows,\Room,\rcut^{*}$ are \tla{ReplayRows},
\tla{ReplayRoom}, \tla{RequiredReplayCut}.

\subsection{Checked invariants and temporal properties}

Table~\ref{tab:props} lists the ten state invariants and ten temporal
properties in the model-checking configurations, with their counterparts in
this paper.

\begin{table}[htbp]
\centering\small
\begin{tabular}{@{}lll@{}}
\toprule
TLA\textsuperscript{+} name & Kind & This paper \\
\midrule
\tla{TypeOK} & invariant & T (Def.~\ref{def:inv}) \\
\tla{LifecycleShape} & invariant & LS \\
\tla{SnapshotDiscipline} & invariant & SD \\
\tla{EmissionDiscipline} & invariant & ED \\
\tla{Capacity} & invariant & CAP \\
\tla{ExactCommittedHistory} & invariant & ECH, eq.~\eqref{eq:ech} \\
\tla{NoPrematureHistory} & invariant & Corollary~\ref{cor:nph} \\
\tla{ScreenCapacity} & invariant & Corollary~\ref{cor:sc};
  Theorem~\ref{thm:capacity} \\
\tla{ReplayShape} & invariant & RS; Lemma~\ref{lem:replay}(i) \\
\tla{NativeSourceSafety} & invariant & NSS \\
\midrule
\tla{HistoryMonotonicity} & action property & Theorem~\ref{thm:actions}(i) \\
\tla{NativeEpochDiscipline} & action property &
  Theorem~\ref{thm:actions}(ii) \\
\tla{FinalImmutability} & action property & Theorem~\ref{thm:actions}(iii) \\
\tla{AppendOnlyMonotonicity} & action property &
  Theorem~\ref{thm:actions}(iv) \\
\tla{ResizeKeepsLogicalHistory} & action property &
  Theorem~\ref{thm:actions}(v) \\
\tla{FailedWriteStops} & invariance & Theorem~\ref{thm:actions}(vi) \\
\tla{StoppedQuiescence} & action property & Theorem~\ref{thm:actions}(vii) \\
\tla{FlushLiveness} & leads-to & Theorem~\ref{thm:flush-live} \\
\tla{ReplayLiveness} & leads-to & Theorem~\ref{thm:replay-live} \\
\tla{QueuedPressureRetirement} & leads-to &
  Theorem~\ref{thm:pressure-live}, Remark~\ref{rem:sharp} \\
\bottomrule
\end{tabular}
\caption{Checked properties and their counterparts.}
\label{tab:props}
\end{table}

\subsection{Configurations and results}

Three boundary configurations were checked exhaustively with TLC
(version 2026.03.31, OpenJDK~17, Apple~M4~Max, 16 worker threads), all with
the snapshot universe
$\Snap=\{\emptyseq,\seq{a},\seq{a,b}\}$ over $\Sigma=\{a,b\}$ (satisfying
Assumption~\ref{ass:params}) and deadlock checking disabled---terminal
states (halted hosts) are intended. No invariant violation, no property
violation, and no liveness counterexample was found in any configuration.

\begin{table}[htbp]
\centering\footnotesize
\begin{tabular}{@{}lccc@{}}
\toprule
 & \tla{ElasticSlots} & \tla{ElasticSlotsBridges}
 & \tla{ElasticSlotsResize} \\
\midrule
$N$ / $H$ & $2$ / $1$ & $1$ / $3$ & $1$ / $1$ \\
$\MaxRes$ / $\MaxLive$ & $1$ / $2$ & $0$ / $1$ & $2$ / $1$ \\
Invariants checked & all $10$ & $7$ & $9$ \\
Temporal properties & all $10$ & $6$ & $8$ \\
States generated & $568{,}654$ & $1{,}736$ & $117{,}042$ \\
Distinct states & $219{,}721$ & $501$ & $52{,}638$ \\
Graph depth & $18$ & $9$ & $16$ \\
Liveness branches & $4$ & $1$ & $1$ \\
Behavior-graph states (liveness) & $878{,}884$ & $501$ & $52{,}638$ \\
Wall time & $22$\,s & $<1$\,s & $3$\,s \\
Fingerprint collision prob.\ & $4.4\times 10^{-9}$
  & $3.4\times 10^{-14}$ & $2.4\times 10^{-10}$ \\
\bottomrule
\end{tabular}
\caption{Exhaustive TLC results (August 24, 2026).}
\label{tab:tlc}
\end{table}

\paragraph{Why these boundaries.}
The main configuration ($N{=}2$, $H{=}1$) maximizes contention: a one-row
transcript forces the overflow/summary path (Theorem~\ref{thm:capacity}(iii))
whenever two blocks are presented, exercises zero-height degradation via
resizes to $h{=}0$, and---with $\MaxLive{=}2$---realizes the persistent count
pressure of Remark~\ref{rem:sharp}, so that
\tla{QueuedPressureRetirement} is checked in exactly the regime where
Theorem~\ref{thm:pressure-live} applies. The bridges configuration
($H{=}3$, no resizes) admits allocations up to $3$ and therefore drives the
deep-shrink trajectory $3\to 2\to 1$ of Lemma~\ref{lem:bridge}(iii) together
with concurrent request/apply interleavings against the reservation
invariant. The resize configuration ($\MaxRes{=}2$) covers
resize-during-replay: a second resize arriving before
\act{ReplaySynchronousSuccess}, which preserves the replay window but resets
preparation, plus all $\{\Preserve,\AppendPol,\Rebuild\}\times$
width/height combinations, epoch resets, and emulator-push capture. Failure
nondeterminism (\act{RetireFailure}, \act{ReplaySynchronousFailure} at every
prefix length up to $\Fmax$) and both exits are enabled in all three
configurations.

\paragraph{Division of labor.}
TLC establishes the theorems of this paper---for the checked
instances---directly on the specification text, including the inductive
stability arguments of Theorem~\ref{thm:safety} whose case analyses are most
error-prone by hand (notably the interaction of RS with the scheduler gate,
and NSS across frontier motion). The hand proofs above are stated for all
parameter values of Assumption~\ref{ass:params} and identify \emph{why} each
invariant is stable, which the model checker does not; where a checked
property is parameter-sensitive we have said so explicitly
(Remark~\ref{rem:sharp}). The two artifacts are kept in lock-step by the
correspondence of Tables~\ref{tab:actions} and~\ref{tab:props}.

% ===========================================================================
\section{Conclusion}\label{sec:conclusion}
% ===========================================================================

Elastic Speculative Slots reconciles speculative, concurrent, elastically
animated live output with the unforgiving append-only nature of terminal
scrollback by decoupling three layers---semantic blocks, an exactly-once
logical ledger, and epoch-scoped physical native history---and by making
every crossing between layers an explicit, provable transition. The
reservation invariant makes concurrent slot animation compositional
(Lemma~\ref{lem:bridge}); the exactness equation $\Ledger=C(c)\cat A(c)$
reduces the entire logical history to two monotone quantities, the frontier
and the head's emitted prefix (Lemma~\ref{lem:head},
Theorem~\ref{thm:safety}); the streaming lemma shows that every abstract
history append is one cursor discipline away from a real VT
(Lemma~\ref{lem:stream}); and the prepared-frame gate turns ``resize replay
without tearing'' from a rendering aspiration into an inductive invariant
(Lemma~\ref{lem:replay}). Failure is fail-stop by construction: a torn write
leaves a tagged, immutable, attributable prefix and a frontier that never
lies. Under weak fairness the protocol drains: replays terminate, flushes
complete, and pressure---when persistent---sweeps finished work into
history exactly once and exactly in order
(Theorems~\ref{thm:replay-live}--\ref{thm:pressure-live}). All of it is
checked, state by state, against \texttt{ElasticSlots.tla}.

\end{document}
